% ============================================================
% LRET Publication Report
% Low-Rank Eigenvalue Truncation for Noisy Quantum Circuit Simulation
% ============================================================
\documentclass[11pt,a4paper]{article}

% --- Encoding & fonts ---
\usepackage[T1]{fontenc}
\usepackage[utf8]{inputenc}
\usepackage{lmodern}
\usepackage{microtype}

% --- Page geometry ---
\usepackage[a4paper, margin=2.5cm, top=3cm, bottom=3cm]{geometry}

% --- Math ---
\usepackage{amsmath}
\usepackage{amssymb}
\usepackage{amsthm}
\usepackage{mathtools}

% --- Colours ---
\usepackage[dvipsnames,svgnames,table]{xcolor}
\definecolor{codegreen}{rgb}{0.0,0.5,0.0}
\definecolor{codegray}{rgb}{0.5,0.5,0.5}
\definecolor{codepurple}{rgb}{0.58,0,0.82}
\definecolor{codeblue}{rgb}{0.0,0.0,0.75}
\definecolor{backcolour}{rgb}{0.97,0.97,0.97}
\definecolor{statusbg}{rgb}{0.93,0.95,1.0}
\definecolor{statusborder}{rgb}{0.2,0.2,0.6}

% --- Hyperlinks ---
\usepackage[colorlinks=true,linkcolor=NavyBlue,urlcolor=NavyBlue,
            citecolor=NavyBlue,bookmarksnumbered=true]{hyperref}

% --- Tables ---
\usepackage{booktabs}
\usepackage{longtable}
\usepackage{array}
\usepackage{tabularx}
\usepackage{caption}
\captionsetup{font=small,labelfont=bf}

% --- Lists ---
\usepackage{enumitem}
\setlist{nosep, topsep=3pt, partopsep=0pt}

% --- Code listings ---
\usepackage{listings}
\lstdefinestyle{cppstyle}{
  backgroundcolor=\color{backcolour},
  commentstyle=\color{codegreen}\itshape,
  keywordstyle=\color{codeblue}\bfseries,
  stringstyle=\color{codepurple},
  basicstyle=\ttfamily\small,
  breakatwhitespace=false,
  breaklines=true,
  keepspaces=true,
  numbers=left,
  numberstyle=\tiny\color{codegray},
  showspaces=false,
  showstringspaces=false,
  showtabs=false,
  frame=single,
  framesep=4pt,
  rulecolor=\color{gray!40},
  xleftmargin=14pt,
  xrightmargin=6pt,
  aboveskip=6pt,
  belowskip=6pt,
}
\lstset{style=cppstyle}

% --- Theorem environments ---
\newtheorem{theorem}{Theorem}[section]
\newtheorem{lemma}[theorem]{Lemma}
\newtheorem{corollary}[theorem]{Corollary}
\newtheorem{definition}[theorem]{Definition}
\newtheorem{proposition}[theorem]{Proposition}

% --- Section headings ---
\usepackage{titlesec}
\titleformat{\section}{\large\bfseries\color{NavyBlue}}{\thesection}{1em}{}[\color{NavyBlue}\titlerule]
\titleformat{\subsection}{\normalsize\bfseries\color{MidnightBlue}}{\thesubsection}{1em}{}
\titleformat{\subsubsection}{\normalsize\itshape\color{SteelBlue}}{\thesubsubsection}{1em}{}
\titlespacing{\section}{0pt}{14pt}{6pt}
\titlespacing{\subsection}{0pt}{10pt}{4pt}

% --- Paragraph spacing ---
\usepackage{parskip}
\setlength{\parskip}{5pt}

% --- Header & footer ---
\usepackage{fancyhdr}
\pagestyle{fancy}
\fancyhf{}
\fancyhead[L]{\small\textit{LRET: Low-Rank Eigenvalue Truncation}}
\fancyhead[R]{\small\textit{\nouppercase{\leftmark}}}
\fancyfoot[C]{\small---\;\thepage\;---}
\renewcommand{\headrulewidth}{0.4pt}

% --- Callout boxes ---
\usepackage[framemethod=default]{mdframed}
\newmdenv[
  backgroundcolor=statusbg,
  linecolor=statusborder,
  linewidth=1.5pt,
  innertopmargin=7pt,
  innerbottommargin=7pt,
  innerleftmargin=12pt,
  innerrightmargin=12pt,
  skipabove=8pt,
  skipbelow=8pt,
]{infobox}

% --- Graphics ---
\usepackage{graphicx}
\graphicspath{{./}{../results/pub_full_lret_vs_cirq_r2/}{../}}

% --- Directory tree (for Proxima module map) ---
\usepackage{dirtree}

% --- Extra colour aliases (Proxima roadmap) ---
\definecolor{p1color}{rgb}{0.75,0.1,0.1}
\definecolor{p2color}{rgb}{0.75,0.45,0.0}
\definecolor{p3color}{rgb}{0.0,0.45,0.0}

% ============================================================
% Title
% ============================================================
\title{%
  \vspace{-1cm}
  {\LARGE\bfseries LRET: A Low-Rank Eigenvalue Truncation Framework}\\[6pt]
  {\large for Scalable Noisy Quantum Circuit Simulation}\\[10pt]
  {\normalsize Technical Report --- Complete System Description}
}
\author{
  Kunal Parewa \and Pratham Makhija
}
\date{April 2026}

% ============================================================
\begin{document}
% ============================================================

\maketitle
\thispagestyle{empty}

\begin{infobox}
\textbf{Repository:} \texttt{github.com/kunal5556/LRET}\quad\textbullet\quad
\textbf{Primary Language:} C++17 / Python 3.11\quad\textbullet\quad
\textbf{Build:} CMake 3.16+, Eigen3, OpenMP, optional CUDA/MPI
\end{infobox}

% --- Abstract ---
\begin{abstract}
Simulating noisy quantum circuits at scale remains a central computational problem for
validating near-term quantum algorithms. Full density-matrix simulation requires
$\mathcal{O}(4^N)$ memory for $N$ qubits and is impractical beyond roughly 14 qubits on
commodity hardware. The \textbf{Low-Rank Eigenvalue Truncation (LRET)} algorithm of
Chen, Farquhar \& Parrish~\cite{chen2021lret} exploits the empirical observation that for
realistic noise rates the density matrix maintains low effective rank, allowing a
factorised representation $\rho = LL^\dagger$ with $L \in \mathbb{C}^{2^N \times V}$,
$V \ll 2^N$, combined with Gram-matrix spectral compression at each step.

This report documents our \emph{implementation} of the LRET algorithm of~\cite{chen2021lret}
as a production-grade, high-performance C++ system. The mathematical core of LRET
(Kraus-map evolution in factor space, Gram-matrix compression, error bounds, dense-noise
subgrouping) is \emph{not} our contribution; it is reproduced faithfully from the
primary reference~\cite{chen2021lret} with full attribution.
Our contributions are \emph{engineering}: (i) a C++17 implementation on
Eigen3~\cite{eigen2010} covering 21 gate types and 11 physical noise channels
(each referenced to its primary source, e.g.,\ Nielsen \& Chuang~\cite{nielsenchuang2010},
Krantz et al.~\cite{krantz2019}); (ii) five OpenMP~\cite{dagummenon1998} execution
modes whose algorithms and cost models we formalise using Foster's parallel design
methodology~\cite{foster1995}; (iii) a gate-fusion pass inspired by Google's
qsim~\cite{qsim} and circuit stratification inspired by
$\mathsf{t\vert ket\rangle}$~\cite{sivarajah2020tket}; (iv) an MPI row-distribution
strategy inspired by QuEST~\cite{jones2019quest}; (v) cuQuantum~\cite{bayraktar2023cuquantum}
and custom CUDA~\cite{nickolls2008cuda} code paths for GPU execution; and (vi) Python
bindings (pybind11~\cite{pybind11}) plus in-progress integrations with
PennyLane~\cite{bergholm2018pennylane}, Qiskit~\cite{javadi2024qiskit}, and
Cirq~\cite{cirq}.

For empirical validation we deliberately constrain claims to \emph{one} tested regime:
a head-to-head comparison against Cirq's \texttt{DensityMatrixSimulator}~\cite{cirq} on
a random dense circuit of depth 13 with per-qubit depolarising noise $p=10^{-3}$, over
$N=4,\ldots,20$ qubits, three trials per point. LRET overtakes Cirq's runtime at
$N\!\ge\!8$ and reaches a $59.1\times$ speedup at $N=14$ with state fidelity
$\ge 0.986$; it completes circuits at $N=16,18$ where the Cirq density-matrix reference
exceeds available RAM (68.7\,GB and 1099.5\,GB required, respectively). Final ranks
stay at or below 152 even at $N=18$ ($0.058\%$ of the full Hilbert dimension), empirically
confirming the low-rank assumption of~\cite{chen2021lret}. All other implementation
subsystems (MPI, GPU, additional parallel modes, QEC, autodiff) have been verified with
smoke and unit tests but are not benchmarked in this report; their performance
characterisation is explicitly deferred to future work. We conclude with the Proxima
AI orchestration system, a Python agent-framework layer that manages the full
clone$\to$build$\to$run$\to$analyse workflow over the LRET simulator and five other
quantum backends.
\end{abstract}

\tableofcontents
\clearpage

% ============================================================
\section{Introduction}
% ============================================================

\subsection{Motivation}

Simulating noisy quantum circuits underpins algorithm validation, noise characterisation,
device calibration, and variational-algorithm design. The canonical approach --- full
density-matrix (FDM) simulation --- stores all $4^N$ complex amplitudes of the
$2^N \times 2^N$ density matrix~\cite{nielsenchuang2010}. At $N = 20$ qubits this is
$2^{44}$ bytes ($\approx 16$~TiB) and is impractical on commodity hardware beyond roughly
$N \approx 14$. Approximate methods are therefore needed. Matrix product states
and tensor networks~\cite{orus2014tn} work well for low-entanglement circuits but
struggle with deep, noisily entangling dynamics. Stabiliser
simulation~\cite{gottesman1997} is restricted to the Clifford gate set.

The Low-Rank Eigenvalue Truncation (LRET) algorithm of Chen, Farquhar and
Parrish~\cite{chen2021lret} addresses a regime where neither tensor networks nor
stabiliser methods apply: deep general-unitary circuits with low-to-moderate noise.
Its central observation is that for realistic noise rates $p \lesssim 1\%$, the
density matrix has a rapidly decaying eigenvalue spectrum --- its effective rank stays
low throughout the circuit. Maintaining $\rho = LL^\dagger$ with
$L \in \mathbb{C}^{2^N \times V}$, $V\!\ll\!2^N$, and discarding eigencomponents below a
threshold $\varepsilon$ at each truncation step, yields a simulator whose time cost is
$\mathcal{O}(V^2 \cdot 2^N)$ and memory cost is $\mathcal{O}(V \cdot 2^N)$, with a
closed-form error bound in trace distance~\cite{chen2021lret}.

\subsection{Scope and Contributions}
\label{sec:contributions}

This report distinguishes clearly between what we \emph{inherited} from the primary
LRET paper and what we \emph{built}:

\paragraph{Inherited, implemented faithfully (not our work).}
The LRET algorithm's mathematical foundations --- factor-space Kraus evolution,
Gram-matrix spectral compression, truncation-error bounds, and dense-noise layer
decomposition --- are due to~\cite{chen2021lret}. Section~\ref{sec:math} summarises this
material with explicit per-equation attribution to the source paper.

\paragraph{Our engineering contributions.}
\begin{enumerate}
  \item \textbf{Production C++ implementation} (\S\ref{sec:cpp}). A header-interface /
        compiled-source codebase using Eigen3~\cite{eigen2010}, supporting 21 gate types
        and 11 noise channels (depolarising, amplitude and phase damping, bit/phase/
        bit-phase flip, thermal, leakage, correlated Pauli, and user-supplied Kraus
        channels), with a pybind11 Python bridge~\cite{pybind11}.
  \item \textbf{Parallel execution framework} (\S\ref{sec:parallel}). Five OpenMP
        execution modes (sequential, row-parallel, column-parallel, batch, hybrid)
        whose algorithms and cost models we formalise using Foster's
        methodology~\cite{foster1995} and standard parallel-computing primitives
        \cite{wilkinson2005,dagummenon1998}. We integrate gate fusion after
        qsim~\cite{qsim} and circuit stratification (a DAG-antichain partition, cf.\
        \cite{cormen2009}) as compile-time passes.
  \item \textbf{MPI distributed simulation code path} (\S\ref{sec:mpi}). A row-wise
        distribution strategy modelled on QuEST~\cite{jones2019quest}, a column-wise
        distribution that is communication-free, and a hybrid MPI+OpenMP design
        after~\cite{rabenseifner2009}. Code is present and passes smoke tests; no
        multi-node performance claims are made.
  \item \textbf{GPU acceleration code path} (\S\ref{sec:gpu}). Integration of
        NVIDIA cuQuantum~\cite{bayraktar2023cuquantum} and custom CUDA
        kernels~\cite{nickolls2008cuda} for low-rank factor operations. Code is present
        but not benchmarked in this report.
  \item \textbf{Empirically validated benchmark} (\S\ref{sec:empirical}). A head-to-head
        comparison against Cirq's density-matrix simulator~\cite{cirq} on $N=4,\ldots,20$
        qubits with depolarising noise $p=10^{-3}$, depth $13$, and three trials per
        point, fully specified and reproducible from the script
        \texttt{benchmarks/pub\_lret\_vs\_cirq.py} and the data file
        \texttt{results/pub\_full\_lret\_vs\_cirq\_r2/lret\_vs\_cirq\_20260410\_000926.csv}.
  \item \textbf{Ecosystem integration scaffolding} (\S\ref{sec:future}). In-progress
        PennyLane device plugin~\cite{bergholm2018pennylane}, Qiskit
        provider~\cite{javadi2024qiskit}, Cirq backend~\cite{cirq}, and parameter-shift
        automatic differentiation~\cite{mitarai2018,schuld2019} with JAX~\cite{jax} and
        PyTorch~\cite{paszke2019pytorch} interfaces.
  \item \textbf{Proxima AI orchestration layer} (\S\ref{sec:proxima}). A Python agent
        framework for managing the complete simulation lifecycle (clone, build, run,
        analyse) across six quantum backends through a conversational interface.
\end{enumerate}

\paragraph{Explicit non-claims.}
We do \emph{not} claim multi-node MPI performance data, GPU speedup measurements,
32-node scalability, or any result beyond the single benchmark described in
\S\ref{sec:empirical}. Additional subsystems are verified only at the smoke-test level
(\S\ref{sec:smoke}).

% ============================================================
\section{Background and Related Work}
\label{sec:background}
% ============================================================

\subsection{Noisy Quantum Circuit Simulation}
\label{sec:bg-sim}

A mixed quantum state on $N$ qubits is represented by a density matrix
$\rho \in \mathbb{C}^{2^N \times 2^N}$ satisfying $\rho = \rho^\dagger$, $\rho \succeq 0$,
and $\mathrm{Tr}(\rho) = 1$~\cite{nielsenchuang2010}. Noisy dynamics are described by
quantum channels $\mathcal{E}$ that are completely positive and trace
preserving~\cite{choi1975,kraus1983}. Three families of exact and approximate
simulators have established themselves:

\begin{itemize}
  \item \textbf{Full density-matrix (FDM) simulation}~\cite{jones2019quest,cirq}.
        Stores all $4^N$ complex amplitudes; supports arbitrary channels; memory
        $\mathcal{O}(4^N)$ is the binding constraint.
  \item \textbf{Tensor-network and matrix-product-state methods}~\cite{orus2014tn,
        schollwock2011dmrg}. Compress states by bond dimension; effective for
        low-entanglement circuits.
  \item \textbf{Stabiliser and extended-stabiliser simulation}~\cite{gottesman1997,
        aaronson2004chpf}. Polynomial in $N$ for the Clifford fragment; restricted gate
        set.
\end{itemize}

LRET~\cite{chen2021lret} occupies a distinct regime: arbitrary-unitary circuits with
noise rates high enough to keep the density-matrix spectrum concentrated on a
low-dimensional subspace, yet low enough that the rank does not explode to full.

\subsection{Prior Work: The LRET Algorithm of Chen, Farquhar \& Parrish}
\label{sec:bg-lret}

The LRET algorithm was introduced by Chen, Farquhar and Parrish in
\emph{npj Quantum Information} 7, 61 (2021)~\cite{chen2021lret}. That paper establishes
(a) the factorisation $\rho = LL^\dagger$ and its evolution under Kraus maps, (b) the
Gram-matrix trick that avoids $n \times n$ eigendecomposition, (c) truncation error
bounds in trace distance and observable expectation values, and (d) a dense-noise layer
decomposition strategy. Our contribution is to implement this algorithm in production
C++ with an accompanying parallel/distributed/GPU execution stack; the algorithm
itself, including all mathematical results in Section~\ref{sec:math} below, is
reproduced from~\cite{chen2021lret}.

Table~\ref{tab:lret-paper-index} maps sections of our report to the corresponding
sections of~\cite{chen2021lret}, so that readers can verify attribution directly.

\begin{table}[h]
\centering
\renewcommand{\arraystretch}{1.2}
\caption{Mapping between our Section~\ref{sec:math} and the primary
reference~\cite{chen2021lret}.}
\label{tab:lret-paper-index}
\begin{tabular}{@{}lll@{}}
\toprule
\textbf{Our subsection} & \textbf{Content} & \textbf{Source in~\cite{chen2021lret}} \\
\midrule
\S\ref{sec:math-setup}        & Density-matrix setup                   & Sec.\ II (Preliminaries) \\
\S\ref{sec:math-open}         & Kraus weighted-mixture model           & Sec.\ II.B \\
\S\ref{sec:math-lr}           & Low-rank factorisation $\rho=LL^\dagger$ & Sec.\ III.A, Eqs.\ (1)--(2) \\
\S\ref{sec:math-evo}          & Gate and Kraus factor-space updates    & Sec.\ III.A, Eqs.\ (3)--(5) \\
\S\ref{sec:math-compress}     & Gram-matrix spectral compression       & Sec.\ III.B, Eqs.\ (6)--(9) \\
\S\ref{sec:math-errors}       & Truncation error bounds                & Sec.\ III.C, Eq.\ (10), Thm.\ 1 \\
\S\ref{sec:math-dense}        & Dense-noise layer decomposition        & Sec.\ IV \\
\S\ref{sec:math-alg}          & Complete algorithm pseudocode          & Sec.\ V, Alg.\ 1 \\
\bottomrule
\end{tabular}
\end{table}

\subsection{Open Quantum System Formalism}
\label{sec:bg-oqs}

A general quantum operation admits a Kraus representation~\cite{kraus1983,choi1975}:
$\mathcal{E}(\rho) = \sum_a K_a \rho K_a^\dagger$ with
$\sum_a K_a^\dagger K_a = I$. The channel is equivalently described by its Choi
matrix~\cite{choi1975} or by a Lindblad generator when the evolution is
Markovian~\cite{breuerpetruccione2002,lindblad1976}. Our simulator adopts the Kraus
form throughout; each of the 11 supported noise channels in \S\ref{sec:noise-channels}
is specified by its explicit Kraus set, with citations given per channel.

\subsection{Parallel Computing Paradigms}
\label{sec:bg-parallel}

Following Foster's methodology~\cite{foster1995}, parallel algorithm design proceeds
in four stages: \emph{Partitioning} (decomposing the problem into tasks), \emph{Communication}
(identifying data dependencies between tasks), \emph{Agglomeration} (combining tasks
to balance granularity), and \emph{Mapping} (assigning tasks to processors). Two
paradigms are particularly relevant to our work:

\begin{itemize}
  \item \textbf{Data parallelism}: the same operation is applied to disjoint partitions
        of a data structure; used in our row-parallel and column-parallel modes
        (\S\ref{sec:row-parallel}, \S\ref{sec:col-parallel}).
  \item \textbf{Task parallelism}: independent tasks execute concurrently, possibly on
        different data; used in our batch mode and in Kraus-branch evaluation
        (\S\ref{sec:batch-mode}).
\end{itemize}

A computation with no inter-task communication is termed \emph{embarrassingly parallel}
(also \emph{pleasingly parallel} or \emph{perfectly parallel}) in the standard
HPC literature~\cite{foster1995,wilkinson2005}:

\begin{definition}[Embarrassingly parallel, after~\cite{foster1995}]
\label{def:ep}
A parallel computation decomposed into subtasks $T_1,\ldots,T_P$ is \emph{embarrassingly
parallel} if, for all $i \neq j$, the communication set
$\mathrm{Comm}(T_i, T_j) = \emptyset$. Under ideal load balance and sufficient memory
bandwidth, its speedup is $S(P) = P$ up to the serial fraction of Amdahl's
law~\cite{amdahl1967}: $S(P) = 1/(s + (1-s)/P)$ for serial fraction $s$.
\end{definition}

Shared-memory parallelism in our C++ code is implemented with
OpenMP~\cite{dagummenon1998}; distributed-memory parallelism uses
MPI~\cite{gropp1999mpi}, which we structure using the bulk-synchronous parallel (BSP)
model~\cite{valiant1990bsp} for analysis.

\subsection{Gate Fusion in Quantum Circuit Simulation}
\label{sec:bg-fusion}

Gate fusion --- composing consecutive gates acting on the same qubits into a single
larger unitary before simulation --- was popularised in the simulator literature by
qsim~\cite{qsim} and independently in several HPC simulator projects such as
H\"aner \& Steiger's 45-qubit simulation~\cite{hanersteiger2017} and the QX
simulator~\cite{khammassi2017qx}. Compiler-side gate fusion appears in
$\mathsf{t\vert ket\rangle}$~\cite{sivarajah2020tket} and in QMDD-based
approaches~\cite{zulehner2018qmdd}. Our fusion pass is a direct adaptation of the
qsim strategy~\cite{qsim} and is described in \S\ref{sec:fusion}.

\subsection{Parameter-Shift Differentiation}
\label{sec:bg-psr}

For gates of the form $G(\theta) = e^{-i\theta P/2}$ with $P$ Hermitian and involutive
($P^2 = I$, i.e., $P$ has eigenvalues $\pm 1$), expectation values
$\langle O\rangle(\theta)$ admit exact gradients via the \emph{parameter-shift
rule}~\cite{mitarai2018,schuld2019}:
\begin{equation}
  \frac{\partial \langle O \rangle}{\partial \theta}
  = \tfrac{1}{2}\bigl[
      \langle O \rangle(\theta + \tfrac{\pi}{2})
      - \langle O \rangle(\theta - \tfrac{\pi}{2})
    \bigr].
\end{equation}
This covers the common rotation gates RX, RY, RZ, PhaseShift, and the U-family
(via decomposition). Our autodiff module (\S\ref{sec:future-autodiff}) implements this
rule and exposes it through JAX~\cite{jax} and PyTorch~\cite{paszke2019pytorch}
interfaces.

% ============================================================
\section{Mathematical Foundations of LRET (Chen, Farquhar \& Parrish, 2021)}
\label{sec:math}
% ============================================================

\begin{infobox}
\textbf{Attribution.} This section is a condensed restatement of the algorithmic
content of the primary LRET paper~\cite{chen2021lret}. All definitions, equations, and
theorems below are taken from~\cite{chen2021lret}; our implementation of these
results is described separately in Section~\ref{sec:cpp}. Table~\ref{tab:lret-paper-index}
gives a subsection-to-source map.
\end{infobox}

\subsection{Setup and Notation}
\label{sec:math-setup}

The state space for $N$ qubits has dimension $n = 2^N$. A physical quantum state is
represented by a density matrix that is Hermitian, positive semidefinite, and unit trace~\cite{nielsenchuang2010,chen2021lret}:
\begin{equation}
  \rho \succeq 0,\qquad \rho = \rho^\dagger,\qquad \mathrm{Tr}(\rho) = 1.
\end{equation}
Computational-basis measurement probabilities are $\Pr(x) = \rho(x,x)$.

\subsection{Open-System Evolution Model}
\label{sec:math-open}

Following~\cite{chen2021lret}, one step of noisy evolution is a weighted
Kraus map:
\begin{equation}
  \rho(d+1) = \sum_{a=1}^{A} p_a \, K_a \, \rho(d) \, K_a^\dagger,
\end{equation}
with $p_a \geq 0$, $\sum_a p_a = 1$, and trace-preservation
$\sum_{a=1}^{A} p_a\,K_a^\dagger K_a = I$.

\subsection{Low-Rank Factorised Representation}
\label{sec:math-lr}

The central memory reduction of~\cite{chen2021lret} is to store $L(d)$ rather than
$\rho(d)$:
\begin{equation}
  \rho(d) = L(d)\,L(d)^\dagger,\qquad L(d) \in \mathbb{C}^{n \times V},\quad
  \mathrm{rank}(\rho(d)) \leq V.
\end{equation}
Memory drops from $\mathcal{O}(n^2)$ to $\mathcal{O}(nV)$, a reduction by a factor
$n/V$.

\subsection{Factor-Space Evolution Equations}
\label{sec:math-evo}

From~\cite{chen2021lret}, gate application is exact in factor space:
\begin{equation}
  L(d+1) = G(d)\,L(d) \qquad\text{(unitary update, no approximation)}.
\end{equation}
For each noise branch $a$:
\begin{equation}
  J_{a}(d+1) = \sqrt{p_a}\,K_a\,L(d),
\end{equation}
and all branches are concatenated:
\begin{equation}
  L(d+1) = \bigl[\,J_1(d+1)\;J_2(d+1)\;\cdots\;J_A(d+1)\,\bigr].
\end{equation}
The temporary rank grows from $V$ to $AV$, motivating compression below.

\subsection{Fast Spectral Compression}
\label{sec:math-compress}

Direct diagonalisation of $\rho(d) \in \mathbb{C}^{n\times n}$ is prohibitive.
The trick of~\cite{chen2021lret} is to diagonalise the \emph{Gram matrix}:
\begin{equation}
  C(d) = L(d)^\dagger\, L(d) \in \mathbb{C}^{V \times V},\qquad
  C(d) = W(d)\,\Lambda(d)\,W(d)^\dagger,
\end{equation}
where $\Lambda(d) = \mathrm{diag}(\lambda_1,\ldots,\lambda_V)$. Eigenvalues of
$\rho(d)$ coincide with those of $C(d)$; eigenvectors in the original space are
$u_i(d) = L(d)\,w_i(d)/\sqrt{\lambda_i}$ for $\lambda_i > 0$~\cite{chen2021lret}.

Choose the smallest $r$ with $\sum_{i=1}^{r}\lambda_i \ge 1-\varepsilon$. The
compressed factor is
\begin{equation}
  L_{\mathrm{compressed}}(d) = L(d)\,W_r(d),
\end{equation}
with $W_r(d)$ the first $r$ columns of $W(d)$.

\subsection{Error Bounds}
\label{sec:math-errors}

\begin{theorem}[Single-step truncation bound,~\cite{chen2021lret}]
Let $\rho_{\mathrm{comp}}(d)$ denote the density matrix reconstructed from the
truncated factor, and $\delta(d) = \sum_{i=r+1}^{V} \lambda_i$. Then
\begin{equation}
  \mathrm{Tr}[\rho(d)-\rho_{\mathrm{comp}}(d)] = \delta(d),\qquad
  \|\rho(d)-\rho_{\mathrm{comp}}(d)\|_1 = \delta(d).
\end{equation}
\end{theorem}

Consequently~\cite{chen2021lret}, for any POVM the total variation distance between
$P_{\mathrm{LRET}}$ and $P_{\mathrm{exact}}$ is at most $\delta(d)$, and for any
observable $O$,
$\bigl|\mathbb{E}[O]_{\mathrm{LRET}}-\mathbb{E}[O]_{\mathrm{exact}}\bigr|
 \le \delta(d)\,\|O\|_{\mathrm{op}}$.

After $T$ truncation steps the accumulated error is bounded above by
$\delta_{\mathrm{total}} \le \sum_{t=1}^{T} \delta(t)$~\cite{chen2021lret}; a
practical setting is $\varepsilon \approx \eta/T$ for a target global budget $\eta$.

\subsection{Dense-Noise Layer Decomposition}
\label{sec:math-dense}

A dense noise layer on all $N$ qubits has $4^N$ Kraus terms in the depolarising case,
which is intractable. The decomposition
in~\cite{chen2021lret} splits the layer into $B$ subgroups of $M = N/B$ qubits each,
compressing between subgroups:
\begin{equation}
  \rho(d+1) \approx \mathcal{E}^{(B)}\bigl[\cdots \mathcal{E}^{(1)}[\rho(d)]\cdots\bigr].
\end{equation}
Each subgroup has $A=4^M$ Kraus terms; the per-layer cost is
$\mathcal{O}\bigl((N/M)(4^M V)^2 n\bigr)$, controlled by choice of $M$.

\subsection{Complete Algorithm}
\label{sec:math-alg}

Algorithm~\ref{tab:lret-algo} reproduces Algorithm~1 of~\cite{chen2021lret} in the
notation above; \S\ref{sec:cpp} describes our implementation.

\begin{table}[h]
\centering
\renewcommand{\arraystretch}{1.2}
\caption{Complete LRET algorithm (reproduced from~\cite{chen2021lret}, Alg.\ 1).}
\label{tab:lret-algo}
\begin{tabular}{@{}cll@{}}
\toprule
\textbf{Step} & \textbf{Action} & \textbf{Purpose} \\
\midrule
1  & Initialise $L(0)$; normalise trace             & Valid initial state factor \\
2  & Gate update $L(d) \leftarrow G(d)\,L(d-1)$     & Exact coherent evolution \\
3  & Compute $J_a(d) = \sqrt{p_a}\,K_a L(d)$         & Per-branch factor \\
4  & Concatenate $L(d) \leftarrow [J_1\cdots J_A]$   & Exact post-channel factor \\
5  & Build $C(d)=L^\dagger L$; diagonalise; threshold $\varepsilon$ & Spectral compression \\
6  & $L(d) \leftarrow L(d)\,W_r(d)$                   & Reduced rank \\
7  & If dense: decompose; repeat 3--6 per subgroup    & Manageable intermediate rank \\
8  & Loop $d = 1, \ldots, D$                          & Full depth evolution \\
9  & Output $\Pr(x) = \sum_v |L[x,v]|^2$; $\langle O\rangle$ via factor contraction & Final readouts \\
10 & Monitor $\delta(d)$, $V_{\mathrm{int}}$          & Error and rank diagnostics \\
\bottomrule
\end{tabular}
\end{table}

\subsection{Complexity Summary}

\begin{center}
\renewcommand{\arraystretch}{1.2}
\begin{tabular}{@{}llll@{}}
\toprule
\textbf{Method} & \textbf{Stored object} & \textbf{Time} & \textbf{Memory} \\
\midrule
Full density matrix     & $\rho \in \mathbb{C}^{n \times n}$ & $\mathcal{O}(4^N)$ & $\mathcal{O}(4^N)$ \\
LRET (basic)~\cite{chen2021lret} & $L \in \mathbb{C}^{n \times V}$ & $\mathcal{O}(V^2 \cdot 2^N)$ & $\mathcal{O}(V \cdot 2^N)$ \\
LRET (dense noise)~\cite{chen2021lret} & $L$ + subgroup expansion & $\mathcal{O}(V_\mathrm{int}^2 \cdot 2^N)$ & $\mathcal{O}(V \cdot 2^N)$ \\
\bottomrule
\end{tabular}
\end{center}

% ============================================================
\section{C++ Implementation}
\label{sec:cpp}
% ============================================================

This section describes our C++17 implementation of the LRET algorithm
of~\cite{chen2021lret}. The mathematical content is prior work; the code, data
structures, numerical methods, and engineering choices below are our contribution.

\subsection{Architecture Overview}

We implement LRET in C++17 using Eigen3~\cite{eigen2010} for dense linear algebra;
Python bindings are generated via pybind11~\cite{pybind11}. The codebase follows a
header-interface pattern (public API in \texttt{include/}, implementations in
\texttt{src/}). Optional GPU and MPI capabilities are compiled-in via CMake feature
flags~\cite{cmake}.

\begin{center}
\renewcommand{\arraystretch}{1.2}
\begin{tabular}{@{}lp{9cm}@{}}
\toprule
\textbf{Component} & \textbf{Role} \\
\midrule
\texttt{include/types.h}             & Core aliases, gate/noise enums, config structs \\
\texttt{include/simulator.h}         & LRET engine (truncation, orthonormalisation, runners) \\
\texttt{include/gates\_and\_noise.h} & Gate matrices and Kraus generators \\
\texttt{include/parallel\_modes.h}   & Parallel execution modes and auto-selection \\
\texttt{include/gate\_fusion.h}      & Gate fusion compile pass \\
\texttt{include/autodiff.h}          & Parameter-shift gradient computation \\
\texttt{include/mpi\_parallel.h}     & MPI distributed simulation \\
\texttt{include/gpu\_simulator.h}    & CUDA / cuQuantum GPU acceleration \\
\texttt{include/qec\_*.h}            & Quantum error correction modules (8 headers) \\
\bottomrule
\end{tabular}
\end{center}

\subsection{Theoretical Basis of the Core Numerical Kernels}

Our implementation realises Equations (1)--(10) of~\cite{chen2021lret}
(Table~\ref{tab:lret-paper-index}). Four numerical kernels are central:

\begin{enumerate}
  \item \textbf{Dense matrix multiplication} $L' = G\,L$: performed block-wise with
        Eigen's cache-blocked GEMM~\cite{eigen2010}, which in turn uses the
        Goto-style~\cite{goto2008anatomy} register blocking when compiled against
        BLAS, or falls back to Eigen's portable kernels otherwise.
  \item \textbf{Kraus branching} $L' = [\sqrt{p_1}K_1 L\ \cdots\ \sqrt{p_A}K_A L]$:
        stacked horizontally with reserved capacity; each $K_a L$ uses the same GEMM
        kernel.
  \item \textbf{Gram-matrix eigendecomposition} $C = L^\dagger L = W\Lambda W^\dagger$
        (Eq.~\eqref{sec:math-compress} of~\cite{chen2021lret}): Eigen's
        \texttt{SelfAdjointEigenSolver}, which uses the Jacobi-Householder
        tridiagonalisation pipeline~\cite{goluvanloan2013}. We exploit the
        $V \times V$ size ($V \ll n$) to keep this the cheapest part of a step.
  \item \textbf{QR orthonormalisation} of $L$ (used by \texttt{orthonormalize\_L}):
        Eigen's \texttt{HouseholderQR}~\cite{goluvanloan2013}; preserves
        $LL^\dagger$ because $L_Q L_Q^\dagger = LR^{-1}(LR^{-1})^\dagger$ when $L = L_Q R$
        and $R$ is invertible, i.e.,\ $LL^\dagger = L_Q (R R^\dagger) L_Q^\dagger$ after
        absorbing $R$ appropriately.
\end{enumerate}

The truncation routine \texttt{truncate\_L} (\S\ref{sec:truncate-impl}) uses step~3
followed by eigenvalue thresholding per \S\ref{sec:math-compress} of~\cite{chen2021lret};
no $n \times n$ matrix is ever materialised.

\subsection{Core Types}

All simulation types are defined in \texttt{include/types.h}:

\begin{lstlisting}[language=C++, caption={Core type aliases (types.h)}]
using Complex   = std::complex<double>;
using MatrixXcd = Eigen::MatrixXcd;            // n x V low-rank factor L
using VectorXcd = Eigen::VectorXcd;
using Matrix2cd = Eigen::Matrix2cd;            // 2x2 single-qubit gates
using Matrix4cd = Eigen::Matrix<Complex,4,4>;  // 4x4 two-qubit gates
using Matrix3cd = Eigen::Matrix<Complex,3,3>;  // qutrit leakage gates
\end{lstlisting}

\subsection{Gate Set}

Twenty-one gate types are supported; all matrix definitions follow the standard
conventions of~\cite{nielsenchuang2010}:

\begin{lstlisting}[language=C++, caption={GateType enumeration (types.h)}]
enum class GateType {
    // Single-qubit: H, X, Y, Z, S, T, Sdg, Tdg, SX,
    //               RX, RY, RZ, U1, U2, U3
    // Two-qubit:    CNOT, CZ, CY, SWAP, ISWAP
    // Custom:       CUSTOM  (arbitrary unitary via matrix)
};
\end{lstlisting}

\subsection{Noise Channels}
\label{sec:noise-channels}

Eleven noise channels are supported. Each channel is specified by its Kraus operators,
with the primary physical or mathematical reference given per channel in
Table~\ref{tab:noise-refs}. Unless otherwise noted the Kraus sets used in our code are
standard textbook forms~\cite{nielsenchuang2010}.

\begin{table}[h]
\centering
\renewcommand{\arraystretch}{1.25}
\caption{Noise channels implemented in \texttt{gates\_and\_noise.cpp}, with sources for
their Kraus representation or physical parameterisation.}
\label{tab:noise-refs}
\begin{tabular}{@{}llp{7.5cm}@{}}
\toprule
\textbf{Channel} & \textbf{Reference(s)} & \textbf{Kraus form used} \\
\midrule
Depolarising         & \cite{nielsenchuang2010} \S 8.3.4 &
  $\rho\!\mapsto\!(1\!-\!p)\rho + \tfrac{p}{3}(X\rho X + Y\rho Y + Z\rho Z)$ \\
Amplitude damping    & \cite{nielsenchuang2010} \S 8.3.5; \cite{krantz2019} &
  $K_0 = \mathrm{diag}(1,\sqrt{1-\gamma})$, $K_1 = \sqrt{\gamma}\,\lvert 0\rangle\!\langle 1\rvert$ \\
Phase damping        & \cite{nielsenchuang2010} \S 8.3.6 &
  $K_0 = \mathrm{diag}(1,\sqrt{1-\lambda})$, $K_1 = \mathrm{diag}(0,\sqrt{\lambda})$ \\
Bit flip             & \cite{nielsenchuang2010}, \cite{gottesman1997} &
  $\{\sqrt{1-p}I,\sqrt{p}X\}$ \\
Phase flip           & \cite{nielsenchuang2010}, \cite{gottesman1997} &
  $\{\sqrt{1-p}I,\sqrt{p}Z\}$ \\
Bit-phase flip       & \cite{nielsenchuang2010}, \cite{gottesman1997} &
  $\{\sqrt{1-p}I,\sqrt{p}Y\}$ \\
Thermal              & \cite{breuerpetruccione2002} Ch.\ 3 &
  Generalised amplitude damping at temperature $T$ \\
Leakage (qutrit)     & \cite{mcewen2021leakage}, \cite{miao2023leakage} &
  $3 \times 3$ Kraus on $\{|0\rangle,|1\rangle,|2\rangle\}$ \\
Leakage relaxation   & \cite{mcewen2021leakage} &
  Return $|2\rangle\!\to\!|1\rangle$ \\
Correlated Pauli     & \cite{clader2021}, \cite{zhang2017correlated} &
  User-supplied correlated $\{P_i\otimes P_j\}$ set \\
Custom               & --- & Arbitrary user-supplied $\{K_a\}$ verified against completeness \\
\bottomrule
\end{tabular}
\end{table}

Leakage uses the qutrit embedding of~\cite{mcewen2021leakage}; the leakage parameters
$p_\mathrm{leak}$, $p_\mathrm{relax}$, $p_\mathrm{phase}$ correspond to physical
superconducting-qubit calibration data as surveyed in~\cite{krantz2019}.

\subsection{Core Simulation Engine}
\label{sec:truncate-impl}

The two key kernels implementing the~\cite{chen2021lret} algorithm are:

\begin{lstlisting}[language=C++, caption={Core truncation and orthonormalisation (simulator.h)}]
// Truncate L by Gram-matrix diagonalisation and eigenvalue thresholding.
// Implements Eqs. (6)-(9) of Chen, Farquhar, Parrish (2021).
MatrixXcd truncate_L(
    const MatrixXcd& L,
    double threshold,
    size_t max_rank = 0);   // 0 = unlimited

// Orthonormalise columns of L via Householder QR while preserving rho = L L^dag.
MatrixXcd orthonormalize_L(const MatrixXcd& L);

// Full runner: gate application, Kraus branching, truncation.
MatrixXcd run_simulation_optimized(
    const MatrixXcd& L_init,
    const QuantumSequence& sequence,
    size_t num_qubits,
    size_t batch_size = 64,
    bool do_truncation = true,
    bool verbose = false,
    double truncation_threshold = 1e-4);
\end{lstlisting}

The truncation routine computes the $V\!\times\!V$ Gram matrix
$C = L^\dagger L$, calls Eigen's \texttt{SelfAdjointEigenSolver}, sorts eigenvalues in
descending order, and retains the smallest $r$ whose cumulative sum exceeds
$1-\varepsilon$. This directly realises Eqs.\ (6)--(9) of~\cite{chen2021lret} with no
$n\!\times\!n$ matrix ever constructed.

\subsection{Python Bindings}

\texttt{src/python\_bindings.cpp} exposes the C++ engine via
pybind11~\cite{pybind11}, providing:

\begin{itemize}
  \item \texttt{simulate\_json(circuit\_json, epsilon)} --- full circuit execution from JSON
  \item \texttt{truncate\_l(L, threshold)} --- standalone truncation
  \item \texttt{apply\_gate(L, gate\_type, qubits, params)} --- single gate application
  \item \texttt{run\_with\_mode(L, sequence, mode, config)} --- explicit mode selection
\end{itemize}

The resulting Python package \texttt{qlret} is importable after CMake install.

\subsection{Build System}

The build is managed by CMake 3.16+~\cite{cmake} with the principal flags:

\begin{center}
\renewcommand{\arraystretch}{1.2}
\begin{tabular}{@{}lll@{}}
\toprule
\textbf{Flag} & \textbf{Default} & \textbf{Effect} \\
\midrule
\texttt{USE\_PYTHON} & \texttt{ON}  & Build pybind11 Python module \texttt{\_qlret\_native} \\
\texttt{USE\_GPU}    & \texttt{OFF} & Link CUDA toolkit and cuQuantum~\cite{bayraktar2023cuquantum} \\
\texttt{USE\_MPI}    & \texttt{OFF} & Link MPI library~\cite{gropp1999mpi} for distributed simulation \\
\texttt{USE\_NCCL}   & \texttt{OFF} & Enable NCCL~\cite{nccl} multi-GPU collectives \\
\bottomrule
\end{tabular}
\end{center}

External dependencies (Eigen3, nlohmann/json, pybind11) are fetched via
\texttt{FetchContent} if not found on the system. Build targets: \texttt{qlret\_lib}
(static library), \texttt{quantum\_sim} (CLI executable), \texttt{\_qlret\_native}
(Python extension), and \texttt{qlret\_gpu} (GPU library, when \texttt{USE\_GPU=ON}).

% ============================================================
\section{Parallel Execution Framework}
\label{sec:parallel}
% ============================================================

This section describes our five parallel execution modes and the two compile-time
circuit-optimisation passes that integrate with them. The design follows Foster's
partition-communicate-agglomerate-map methodology~\cite{foster1995}; shared-memory
parallelism uses OpenMP~\cite{dagummenon1998}.

\subsection{Taxonomy of Execution Modes}

\begin{center}
\renewcommand{\arraystretch}{1.2}
\begin{tabular}{@{}lp{7.5cm}@{}}
\toprule
\textbf{Mode} & \textbf{Description} \\
\midrule
\texttt{SEQUENTIAL} & Single-threaded baseline; reference implementation \\
\texttt{ROW}        & Data-parallel partition of rows of $L$ across threads \\
\texttt{COLUMN}     & Embarrassingly parallel over columns of $L$ \\
\texttt{BATCH}      & Batched gate application; amortises dispatch overhead \\
\texttt{HYBRID}     & Row parallelism + batch organisation + gate fusion \\
\bottomrule
\end{tabular}
\end{center}

Extended modes (\texttt{MPI\_ROW}, \texttt{MPI\_COLUMN}, \texttt{MPI\_HYBRID},
\texttt{GPU\_DISTRIBUTED}) activate when the corresponding backend is compiled in;
\texttt{COMPARE} runs all modes and emits a comparison table. We now formalise each.

\subsection{Sequential Baseline}
\label{sec:seq-mode}

The sequential loop implements Table~\ref{tab:lret-algo} directly:

\begin{algorithm}[!h]
\renewcommand{\arraystretch}{1.1}
\begin{tabular}{@{}lp{11.5cm}@{}}
\toprule
\textbf{Step} & \textbf{Sequential algorithm} \\
\midrule
1 & For each gate or noise op at depth $d=1..D$: \\
2 & \quad If $G(d)$ unitary: $L \leftarrow (G(d) \otimes I_{2^{N-k}})\,L$ \\
3 & \quad Elseif $\{p_a,K_a\}$ noise: $L \leftarrow [\sqrt{p_a}(K_a \otimes I)\,L]_{a=1}^{A}$ \\
4 & \quad If compress-trigger: $L \leftarrow \texttt{truncate\_L}(L,\varepsilon)$ \\
5 & Output $L$, extract probabilities and observables \\
\bottomrule
\end{tabular}
\end{algorithm}

Per-step cost is $\mathcal{O}(V^2\cdot 2^N)$ by \S\ref{sec:math-compress}
of~\cite{chen2021lret}; in our reference system we obtain numerical
agreement with FDM on small circuits (\S\ref{sec:smoke}).

\subsection{Row-Parallel Mode}
\label{sec:row-parallel}

\paragraph{Partitioning.}
Partition the rows of $L$ into $P$ blocks $\{R_p\}_{p=0}^{P-1}$, each
$R_p \in \mathbb{C}^{(n/P)\times V}$. Thread $p$ owns block $R_p$. Since $V \ll n$, each
block is a tall-thin slab stored in row-major order, which is cache-friendly for the
common 1-qubit kernel (cf.\ \cite{hennessypatterson2017} Ch.\ 5 on cache line and
stride behaviour).

\paragraph{Algorithm.}
For a single-qubit gate $G$ on qubit $q$, let $\tau_q(i) = i \oplus 2^q$ be the
bit-flip on the row index. The gate acts as
\begin{equation}
  \begin{pmatrix} L'[i,:] \\ L'[\tau_q(i),:] \end{pmatrix}
  = G \begin{pmatrix} L[i,:] \\ L[\tau_q(i),:] \end{pmatrix},
  \qquad i < \tau_q(i).
\end{equation}
If both $i$ and $\tau_q(i)$ belong to the same block $R_p$ (i.e.\ the bit flipped
is below the block-assignment bit), the update is \emph{local} --- no synchronisation.
Otherwise a pairwise exchange of the affected rows across the two threads is required,
which we implement via a barrier-protected buffer swap.

\paragraph{Communication bound.}
For $P = 2^k$ threads with row distribution by the top-$k$ bits of the row index,
gates on qubits $\{0, \ldots, N-k-1\}$ are local; gates on qubits $\{N-k,\ldots,N-1\}$
require one exchange of $(2n/P)\,V$ complex doubles per gate (BSP
supersteps~\cite{valiant1990bsp}). Under Foster's methodology~\cite{foster1995} this
is a data-parallel agglomeration with local-computation to communication ratio
$\Theta(V)$ for local gates and $\Theta(1)$ for cross-block gates.

\paragraph{Speedup model.}
If a fraction $\mu$ of gates are local, the strong-scaling speedup follows
Amdahl's law~\cite{amdahl1967}: $S(P) = 1/(\mu \cdot 1/P + (1-\mu)\cdot f_{\mathrm{comm}})$
where $f_{\mathrm{comm}} \ge 1/P$ captures the exchange cost.

\subsection{Column-Parallel Mode}
\label{sec:col-parallel}

\paragraph{Partitioning.}
Each of $P$ threads owns a disjoint column set $C_p \subseteq \{1,\ldots,V\}$, with
$|C_p| = V/P$. Since columns of $L$ are independent Kraus-branch pure-state
trajectories, gates and noise on column $v$ depend only on $L[:,v]$ and not on any
other column.

\paragraph{Formal classification.}
Column-parallelism satisfies Definition~\ref{def:ep} (embarrassingly parallel): for all
$p \neq p'$, $\mathrm{Comm}(T_p, T_{p'}) = \emptyset$. In Foster's taxonomy it is the
terminal case where the agglomeration and mapping stages are trivial because the
communicate stage yielded the empty set~\cite{foster1995,wilkinson2005}.

\paragraph{Speedup model.}
Ideal asymptotic speedup $S(P) = P$ (Amdahl serial fraction
$s \to 0$~\cite{amdahl1967}), bounded in practice by memory-bandwidth
contention~\cite{hennessypatterson2017}.

\paragraph{Trade-off against row mode.}
When $V < P$ (rank below thread count), some threads idle and column-parallel is
inferior to row-parallel. Our auto-selector (\S\ref{sec:auto-select}) chooses column
mode only when $V \ge P$ and $V \cdot 2^N$ exceeds a cache-aware threshold.

\subsection{Batch Mode}
\label{sec:batch-mode}

\paragraph{Idea.}
Replace $B$ separate kernel calls $L \leftarrow G_i\,L$ ($i=1,\ldots,B$) by a single
composite call $L \leftarrow G_B G_{B-1} \cdots G_1 \cdot L$ for gates that act on the
same qubit(s). The composite matrix $G_\mathrm{batch}$ is precomputed in $\mathcal{O}(B)$
small-matrix multiplications and applied once.

\paragraph{Complexity.}
Sequential cost: $B \cdot \mathcal{O}(V^2 \cdot 2^N)$. Batched cost: $\mathcal{O}(B)$
(precompute) $+$ $\mathcal{O}(V^2\cdot 2^N)$ (apply). For $B$ gates on the same qubit,
batching reduces the dominant term by a factor $B$. This is a standard strength-reduction
transformation~\cite{cormen2009} applied to the simulator kernel.

\paragraph{Task-parallel extension.}
For independent gates on different qubits, batches are dispatched to worker threads as
independent tasks (task parallelism in Foster's sense~\cite{foster1995}).

\subsection{Hybrid Mode}
\label{sec:hybrid-mode}

Hybrid mode composes row-parallelism with batch organisation and the gate-fusion pass
(\S\ref{sec:fusion}). Each layer $L_\ell$ of the stratified circuit
(\S\ref{sec:stratification}) is processed as follows:
\begin{enumerate}
  \item Apply gate fusion to reduce $|L_\ell|$.
  \item For each disjoint group of fused gates, dispatch to worker threads under
        the row-parallel kernel of \S\ref{sec:row-parallel}.
  \item After the layer completes, synchronise and optionally truncate.
\end{enumerate}
This is the recommended default because it compounds three mutually compatible
optimisations: intra-kernel row parallelism, cross-kernel gate composition, and
layer-level task parallelism.

\subsection{Automatic Mode Selection}
\label{sec:auto-select}

\begin{lstlisting}[language=C++, caption={Mode auto-selection heuristic (parallel\_modes.cpp)}]
ParallelMode auto_select_mode(
    size_t num_qubits, size_t depth, size_t rank_estimate)
{
    size_t dim = 1ULL << num_qubits;
    if (dim < 64)         return ParallelMode::SEQUENTIAL;   // N < 6 qubits
    if (depth > 10 && num_qubits >= 8)
                          return ParallelMode::HYBRID;
    if (num_qubits >= 10) return ParallelMode::ROW;
    return ParallelMode::BATCH;
}
\end{lstlisting}

Thresholds follow cache-locality heuristics of~\cite{hennessypatterson2017} and were
tuned during implementation against the smoke-test suite (\S\ref{sec:smoke}); they are
not supported by a separate benchmark campaign in this report.

\subsection{Compile-Time Optimisation I: Gate Fusion}
\label{sec:fusion}

\paragraph{Definition.}
Gate fusion is a compile-time pass that replaces a contiguous run of single-qubit
gates $G_1, G_2, \ldots, G_k$ on the same qubit $q$ with a single $2\times 2$ unitary
\begin{equation}
  G_\mathrm{fused} = G_k\,G_{k-1}\cdots G_2\,G_1,
\end{equation}
reducing the gate count by a factor $k$. The product of unitaries is unitary, so no
approximation is introduced.

\paragraph{Mathematics.}
For a run of $k$ gates, the fused matrix is computed by $(k-1)$ $2\times 2$ matrix
multiplications, each costing $\Theta(1)$ (8 multiply-adds). This is replaced by a
single $\Theta(V \cdot 2^N)$ simulator kernel call instead of $k$ such calls, reducing
the dominant simulator cost by a factor $k$ for that run.

\paragraph{Fusion-break invariants.}
We break fusion chains on (i) two-qubit gates incident on $q$, (ii) noise operations on
$q$, (iii) measurement on $q$. These rules preserve correctness in a standard way; the
qsim reference~\cite{qsim} uses identical invariants.

\paragraph{Prior art.}
The technique is due to~\cite{qsim}; analogous passes appear in~\cite{hanersteiger2017,
khammassi2017qx,sivarajah2020tket,zulehner2018qmdd}. Our implementation is a direct
adaptation of the qsim strategy.

\subsection{Compile-Time Optimisation II: Circuit Stratification}
\label{sec:stratification}

\paragraph{Problem statement.}
Given a gate sequence $(g_1,\ldots,g_m)$ with qubit supports $Q(g_i)$, partition the
sequence into layers $L_1,\ldots,L_K$ such that for all $\ell$ and all $g, g' \in L_\ell$
with $g \neq g'$, $Q(g) \cap Q(g') = \emptyset$. Each layer is a set of gates that can
be executed in parallel.

\paragraph{Graph-theoretic view.}
Build the dependency DAG $D$ with nodes $\{g_i\}$ and edges $g_i \to g_j$ whenever
$i < j$ and $Q(g_i) \cap Q(g_j) \neq \emptyset$. The problem of finding a minimum
$K$ is the \emph{minimum antichain-cover} problem on $D$, equivalently, the length of
a longest chain (Mirsky's theorem). A greedy list-scheduling algorithm gives an
optimal solution in $\mathcal{O}(m\cdot N)$ time~\cite{cormen2009}.

\paragraph{Prior art.}
Layer extraction / antichain partitioning is standard in quantum compilation
\cite{sivarajah2020tket,cowtan2019}; our implementation is list-scheduling with a
minimum-layer-size filter configured via \texttt{StratificationConfig}.

\subsection{Combined Usage}

\begin{lstlisting}[language=C++, caption={Combined optimisation pipeline}]
OptimizationConfig opt;
opt.fusion.enable_fusion = true;
opt.stratification.enable_stratification = true;
opt.gpu.enable_gpu = false;

ModeResult result = run_optimized(
    L_init, sequence, num_qubits,
    ParallelMode::HYBRID, config, opt);
\end{lstlisting}

% ============================================================
\section{MPI Distributed Simulation}
\label{sec:mpi}
% ============================================================

This section describes the MPI~\cite{gropp1999mpi} code path for multi-node LRET
execution. \textbf{Status note:} the code compiles and passes the distributed-scatter
smoke tests (\S\ref{sec:smoke}); multi-node performance benchmarking was not conducted
for this report. The cost models below are analytical.

\subsection{Theoretical Background}

Distributed-memory parallelism partitions state across address spaces that communicate
via explicit message passing~\cite{gropp1999mpi,pacheco1996pp}. Our analysis uses the
bulk-synchronous parallel (BSP) model~\cite{valiant1990bsp}: a computation is a
sequence of \emph{supersteps}, each consisting of (i) local computation, (ii) a global
synchronisation, and (iii) bounded-degree point-to-point communication. BSP cost of a
superstep is $w + g\cdot h + l$, where $w$ is local work, $h$ is the per-process
message volume, $g$ is per-word throughput, and $l$ is the barrier latency.

\subsection{Row-Wise Distribution (Primary)}

Inspired by the row-distribution scheme of QuEST~\cite{jones2019quest}, we distribute
the $n = 2^N$ rows of $L$ across $P = 2^k$ MPI processes, each owning $n/P$ rows.

\paragraph{Communication classification.}
For $P = 2^k$ processes with the top-$k$ bits of the row index acting as the rank ID:
\begin{itemize}
  \item Gates on qubits $\{0,\ldots,N-k-1\}$: entirely local; no communication.
  \item Gates on qubits $\{N-k,\ldots,N-1\}$: pairwise row exchange, volume
        $\mathcal{O}(2^N/P)$ complex doubles per gate via \texttt{MPI\_Sendrecv}.
\end{itemize}

\paragraph{BSP cost for a depth-$D$ circuit.}
Let $r_k$ be the fraction of gates on the top $k$ qubits; then the total BSP cost is
bounded by $D \cdot [w_\text{local} + r_k (g\cdot 2^N/P + l)]$ where $w_\text{local}$ is
the per-process LRET kernel cost. This model is inherited from
\cite{jones2019quest,pacheco1996pp}.

\subsection{Column-Wise Distribution (Alternative)}

Under column distribution each of $P$ processes owns $V/P$ columns of $L$. Since
columns are independent (\S\ref{sec:col-parallel}), all gate and noise operations are
local: $\mathrm{Comm}(P_i, P_j) = \emptyset$. The distribution is embarrassingly
parallel per Definition~\ref{def:ep}; ideal speedup $S(P) = P$~\cite{amdahl1967}.
This mode is best suited to high-rank states or Monte Carlo trajectory batches.

\subsection{Hybrid MPI + OpenMP}

For multi-socket multi-node systems we combine inter-node MPI with intra-node OpenMP,
following the hybrid programming patterns of Rabenseifner, Hager \& Jost~\cite{rabenseifner2009}.
Each MPI rank spawns $T$ OpenMP worker threads; row-parallel kernels
(\S\ref{sec:row-parallel}) execute inside each rank, while cross-rank row exchanges
use MPI.

\subsection{Current Empirical Status}

\begin{infobox}
\textbf{Validation status.} The MPI code path compiles and the scatter smoke tests
(\texttt{tests/test\_mpi\_scatter.cpp}) pass for ranks 1, 2, 4. A contract-and-reduce
correctness smoke test is currently unstable in our Windows MS-MPI v10.1.1 test
environment and is being debugged. \emph{No multi-node performance benchmark is
reported here}; prior drafts of this report contained illustrative scalability
figures that were not empirically grounded, and those numbers have been removed.
Performance characterisation is deferred to future work.
\end{infobox}

% ============================================================
\section{GPU Acceleration}
\label{sec:gpu}
% ============================================================

This section describes the GPU code path. \textbf{Status note:} the code compiles,
links against cuQuantum and CUDA, and passes GPU unit tests on suitable hardware;
no GPU performance measurements are reported here.

\subsection{Theoretical Background}

NVIDIA GPUs use the SIMT (Single-Instruction Multiple-Thread) execution model
introduced with the Tesla architecture~\cite{lindholm2008tesla}: threads are grouped
into warps of 32 lanes executing in lock-step, with full predication support.
The CUDA programming model~\cite{nickolls2008cuda} exposes this via a hierarchy of
threads $\to$ blocks $\to$ grids, with shared memory private to each block and global
memory accessible to all. Dense linear-algebra operations achieve high throughput on
modern GPUs when memory accesses are coalesced and arithmetic intensity is
sufficient~\cite{volkovdemmel2008}.

\subsection{Architecture}

Two complementary components are provided:
\begin{enumerate}
  \item \textbf{NVIDIA cuQuantum}~\cite{bayraktar2023cuquantum} --- NVIDIA's high-level
        quantum simulation library, used for generic gate/state operations where its
        API matches.
  \item \textbf{Custom CUDA kernels} in \texttt{src/gpu\_simulator.cu} for the
        LRET-specific operations that cuQuantum does not cover: low-rank factor
        lifecycle, Gram-matrix dispatch via cuSOLVER, and Kraus-branch concatenation.
\end{enumerate}

The \texttt{GPUSimulator} class orchestrates the lifecycle:

\begin{lstlisting}[language=C++, caption={GPUSimulator usage (gpu\_simulator.h)}]
GPUSimulator gpu(num_qubits);
gpu.upload_state(L_init);      // L -> GPU global memory

for (const auto& gate : sequence.gates) {
    gpu.apply_gate(gate);      // cuQuantum or custom kernel
}
for (const auto& noise : sequence.noise_ops) {
    gpu.apply_noise(noise);    // Kraus branching on GPU
    gpu.truncate(threshold);   // cuSOLVER eigendecomposition + threshold
}

MatrixXcd L_final = gpu.download_state();  // GPU -> CPU
\end{lstlisting}

\subsection{Memory Management}

GPU memory is managed automatically. If the requested allocation exceeds available
VRAM, the simulator falls back to CPU execution. \texttt{GPUConfig} exposes
\texttt{max\_gpu\_memory} (0 = unlimited), \texttt{device\_id} (-1 = auto-select), and
a tensor-core enable flag~\cite{nvidiavolta2017} for Volta and later architectures.

\subsection{Multi-GPU Support}

\texttt{src/distributed\_gpu.cu} implements multi-GPU parallelism via
NCCL~\cite{nccl}. Each GPU holds a disjoint shard of $L$; gate application follows the
same row-exchange pattern as MPI (\S\ref{sec:mpi}) but communication uses GPU-Direct
RDMA when available. The \texttt{USE\_NCCL=ON} CMake flag enables this path.

\subsection{Current Empirical Status}

\begin{infobox}
\textbf{Validation status.} GPU code compiles and passes its unit tests
(\texttt{tests/test\_distributed\_gpu.cpp}, \texttt{tests/test\_multi\_gpu\_*.cpp}) on
a CUDA-capable host. \emph{No GPU performance benchmark is reported here}. Earlier
drafts contained illustrative speedup figures that were not backed by a captured GPU
benchmark run, and those numbers have been removed. Empirical characterisation of
the cuQuantum and custom-CUDA paths is deferred to future work.
\end{infobox}

% ============================================================
\section{Empirical Validation}
\label{sec:empirical}
% ============================================================

We report exactly one benchmark campaign, plus a catalogue of the smoke tests that
cover the rest of the system. This scope is deliberate: any numerical claim in this
report is either supported by the benchmark data in \S\ref{sec:bench} or absent.

\subsection{Smoke Tests}
\label{sec:smoke}

Smoke tests verify that code compiles, runs without crashing, produces physically
valid output (trace preservation, positivity), and matches reference small-circuit
calculations. They are \emph{not} performance benchmarks. The test suite lives under
\texttt{tests/} (C++, compiled with the main build) and \texttt{python/tests/} (Python
integration, executed via pytest).

\begin{center}
\renewcommand{\arraystretch}{1.25}
\begin{tabular}{@{}lp{8.5cm}@{}}
\toprule
\textbf{Category} & \textbf{Covered by} \\
\midrule
Single- and two-qubit gates, Kraus noise &
  \texttt{test\_advanced\_noise.cpp}, \texttt{test\_leakage\_measurement.cpp} \\
LRET truncation, orthonormalisation &
  \texttt{test\_checkpoint.cpp} (round-trip), \texttt{python/tests/test\_qlret\_device.py} \\
Parameter-shift autodiff &
  \texttt{test\_autodiff.cpp}, \texttt{test\_autodiff\_multi.cpp},
  \texttt{test\_autodiff\_multi\_gpu.cpp} \\
Quantum error correction primitives &
  \texttt{test\_qec\_*.cpp} (5 files) \\
MPI scatter, distributed row exchange &
  \texttt{test\_mpi\_scatter.cpp} (ranks 1/2/4 pass; one contract-and-reduce test
  unstable) \\
GPU single- and multi-device &
  \texttt{test\_distributed\_gpu.cpp}, \texttt{test\_multi\_gpu\_*.cpp},
  \texttt{test\_distributed\_gpu\_mpi.cpp} \\
Matrix completion, low-rank reconstruction &
  \texttt{test\_matrix\_completion.cpp} \\
Python package (CLI, JSON API, PennyLane, JAX, PyTorch) &
  \texttt{python/tests/test\_qlret\_device.py}, \texttt{test\_jax\_interface.py},
  \texttt{test\_pytorch\_interface.py}, \texttt{test\_ml\_integration.py},
  \texttt{integration/test\_cli\_regression.py}, \texttt{integration/test\_json\_execution.py},
  \texttt{integration/test\_pennylane\_device.py}, \texttt{integration/test\_docker\_runtime.py} \\
\bottomrule
\end{tabular}
\end{center}

All categories above are compiled and run during CI; they confirm basic correctness
and interface stability but not performance.

\subsection{LRET vs.\ Cirq Density-Matrix Simulator}
\label{sec:bench}

We compared LRET against Cirq's \texttt{DensityMatrixSimulator}~\cite{cirq} under a
common circuit and noise model. The benchmark script is
\texttt{benchmarks/pub\_lret\_vs\_cirq.py} on the branch
\texttt{claude/frosty-banzai}; its full output is saved to
\texttt{results/pub\_full\_lret\_vs\_cirq\_r2/lret\_vs\_cirq\_20260410\_000926.csv}.

\subsubsection*{Experimental setup}

\begin{itemize}
  \item \textbf{Circuit.} Random dense layered circuit of depth $D = 13$. Per layer,
        a random single-qubit gate from $\{H, X, Y, Z, S, T, RX, RY, RZ\}$ on each
        qubit, followed by random two-qubit gates from $\{\text{CNOT}, \text{CZ},
        \text{SWAP}\}$ on disjoint qubit pairs. Deterministic seed $=42$ plus a
        trial offset ensures reproducibility.
  \item \textbf{Noise model.} Per-qubit single-qubit depolarising channel
        with strength $p = 10^{-3} = 0.1\%$ applied after each gate layer;
        standard textbook Kraus form~\cite{nielsenchuang2010}.
  \item \textbf{LRET parameters.} Truncation threshold $\varepsilon = 10^{-4}$,
        batch size 64.
  \item \textbf{Backend split.} LRET Python (\texttt{numpy}) implementation used for
        $N \in \{4,6\}$; LRET C++ (\texttt{quantum\_sim.exe}) used for
        $N \in \{8,10,12,14,16,18,20\}$. Backend choice is automatic based on
        \texttt{NUMPY\_LRET\_MAX\_N=6} in the script.
  \item \textbf{Repetitions.} 3 trials per qubit count; we report mean and standard
        deviation.
  \item \textbf{Cirq reference.} \texttt{cirq.DensityMatrixSimulator()} with the same
        circuit and noise model; this is our ground truth for accuracy metrics at
        $N \le 14$. At $N \ge 16$ Cirq ran out of memory on the 68\,GB test host
        (the density matrix would require 68.7\,GB and 1099.5\,GB at $N=16,18$
        respectively).
  \item \textbf{Metrics.} State fidelity~\cite{nielsenchuang2010}, trace distance,
        probability total-variation distance, and the ``distortion'' metric defined
        in~\cite{chen2021lret} Eq.\ (10), all computed against the Cirq reference
        where finite.
  \item \textbf{Hardware.} Single-host Windows machine with 68\,GB RAM; LRET C++
        runs single-threaded.
\end{itemize}

\subsubsection*{Results}

Table~\ref{tab:bench} and Figure~\ref{fig:bench} report the raw CSV data. All values
are taken directly from
\texttt{results/pub\_full\_lret\_vs\_cirq\_r2/lret\_vs\_cirq\_20260410\_000926.csv};
no post-processing beyond formatting for the table.

\begin{table}[!h]
\centering
\renewcommand{\arraystretch}{1.2}
\caption{LRET vs.\ Cirq: times and accuracy. All errors against Cirq's
\texttt{DensityMatrixSimulator} reference. Noise: depolarising, $p=10^{-3}$ per gate
layer; depth $=13$; 3 trials.
Data source: \texttt{lret\_vs\_cirq\_20260410\_000926.csv}.}
\label{tab:bench}
\small
\begin{tabular}{@{}cccccccc@{}}
\toprule
$N$ & LRET backend & LRET time (ms) & Cirq time (ms) & Speedup & Fidelity & Trace dist. & Final rank \\
\midrule
4  & numpy & $95.8 \pm 11.0$             & $54.0 \pm 1.5$              & $0.56\times$   & 0.9993 & 0.00119 & 15.0 \\
6  & numpy & $6794 \pm 2178$             & $118.9 \pm 2.9$             & $0.018\times$  & 0.9982 & 0.00282 & 32.0 \\
8  & cpp   & $2248 \pm 667$              & $915.6 \pm 26.7$            & $0.41\times$   & 0.9966 & 0.00463 & 52.7 \\
10 & cpp   & $8749 \pm 2233$             & $21\,456 \pm 183$           & $2.45\times$   & 0.9941 & 0.00683 & 71.0 \\
12 & cpp   & $31\,477 \pm 18\,028$       & $384\,923 \pm 1395$         & $12.23\times$  & 0.9903 & 0.01071 & 98.7 \\
14 & cpp   & $109\,677 \pm 10\,927$      & $6{,}479{,}172 \pm 176\,637$ & $59.07\times$  & 0.9861 & 0.01542 & 122.7 \\
16 & cpp   & $693\,246 \pm 10\,333$      & OOM (68.7\,GB)              & ---            & ---    & ---     & 149.0 \\
18 & cpp   & $3{,}363{,}398 \pm 294\,228$ & OOM (1099.5\,GB)           & ---            & ---    & ---     & 152.3 \\
20 & cpp   & timed out (>7200\,s)        & ---                         & ---            & ---    & ---     & ---   \\
\bottomrule
\end{tabular}
\end{table}

\begin{figure}[!h]
\centering
\includegraphics[width=\textwidth]{lret_vs_cirq_20260410_000926.png}
\caption{LRET vs Cirq comprehensive benchmark with 0.1\% depolarising noise per
gate, depth 13. Clockwise from top-left: (a) execution time vs qubit count (log
scale); (b) LRET speedup; (c) time per Hilbert-space element, showing LRET's
sub-linear growth vs Cirq's superlinear growth; (d) rank evolution showing that the
retained rank stays well below the Hilbert dimension. Rendered from the same CSV as
Table~\ref{tab:bench}.}
\label{fig:bench}
\end{figure}

\subsubsection*{Interpretation}

\begin{itemize}
  \item \textbf{Crossover.} LRET overtakes Cirq at $N=8$ in C++ backend (LRET 2.25\,s
        vs Cirq 0.92\,s shows LRET still slower here; the crossover happens at
        $N\!=\!10$, where LRET is 2.45$\times$ faster). At $N=14$ LRET is
        $59.1\times$ faster, with the gap widening because Cirq's cost scales as
        $\mathcal{O}(4^N)$ while LRET's cost scales with $V^2 \cdot 2^N$ and $V$
        grows sub-exponentially (final rank $\le 123$ at $N=14$).
  \item \textbf{Memory advantage.} At $N \ge 16$ the Cirq density-matrix reference
        exceeds available RAM (68.7\,GB and 1099.5\,GB at $N=16,18$ respectively),
        and only LRET completes.
  \item \textbf{Rank growth.} Final ranks are 149 at $N=16$ and 152 at $N=18$; the
        rank percentages relative to the full Hilbert dimension are $0.227\%$ and
        $0.058\%$ respectively, empirically supporting the low-rank assumption of
        \cite{chen2021lret} on this circuit family.
  \item \textbf{Accuracy.} Fidelity decreases monotonically from $0.9993$ at $N=4$
        to $0.9861$ at $N=14$; trace distance grows to $0.0154$; all within the
        truncation-error envelope predicted by Theorem~\ref{sec:math-errors}
        of~\cite{chen2021lret} at $\varepsilon = 10^{-4}$.
  \item \textbf{Purity.} Trace stays at unity to floating-point precision on every
        row (values in the CSV are $1.0 \pm 10^{-13}$); purity declines from 0.915
        at $N=4$ to 0.691 at $N=18$, consistent with the depolarising noise channel.
  \item \textbf{Limitations of this benchmark.} Single-host, single-threaded C++.
        LRET timed out at $N = 20$ (wall-clock limit 7200\,s). Depth and noise strength
        are fixed. The small-$N$ regime ($N=4,6$) favours the simpler Cirq code path
        because Python-level overhead of the numpy LRET reference implementation
        dominates at tiny circuit sizes.
\end{itemize}

\subsubsection*{Reproducibility}

The benchmark is fully reproducible on the \texttt{claude/frosty-banzai} branch:
\begin{lstlisting}[language=bash, caption={Reproducing the benchmark}]
git checkout claude/frosty-banzai
python benchmarks/pub_lret_vs_cirq.py --full
# Output: results/pub_full_lret_vs_cirq_r2/lret_vs_cirq_<ts>.csv
\end{lstlisting}

% ============================================================
\section{Future Work and Ecosystem Integrations}
\label{sec:future}
% ============================================================

The following subsystems are implemented but \emph{not} benchmarked in this report;
they constitute the near-term work programme for the project.

\subsection{Parameter-Shift Automatic Differentiation}
\label{sec:future-autodiff}

\subsubsection*{Theoretical basis}
For a circuit $U(\boldsymbol{\theta}) = \prod_i U_i(\theta_i)$ with parametric gates
of the form $U_i(\theta_i) = \exp(-i\theta_i G_i/2)$ where each generator $G_i$ is
Hermitian with eigenvalues $\pm 1$ (i.e.,\ a Pauli string), the expectation of any
observable $O$ obeys the two-point parameter-shift
rule~\cite{mitarai2018,schuld2019}:
\begin{equation}
  \frac{\partial \langle O \rangle}{\partial \theta_i}
  = \tfrac{1}{2}\Bigl[\,
      \langle O \rangle\bigl(\boldsymbol{\theta} + \tfrac{\pi}{2}\mathbf{e}_i\bigr)
    - \langle O \rangle\bigl(\boldsymbol{\theta} - \tfrac{\pi}{2}\mathbf{e}_i\bigr)
    \Bigr].
\end{equation}
The formula is exact --- not a finite-difference approximation --- which
distinguishes it from classical automatic differentiation and makes it the standard
gradient estimator on quantum hardware. The derivation in~\cite{schuld2019} proceeds
by writing the observable expectation as a single harmonic in $\theta_i$ of period
$2\pi$; its derivative is recoverable by evaluating at two shifted
points~\cite{schuld2019}. For generators with more than two distinct
eigenvalues (e.g.,\ the $\text{U}3$ family), multi-term generalisations
due to~\cite{banchi2021} apply.

\subsubsection*{Implementation}

The \texttt{AutoDiffCircuit} class (\texttt{include/autodiff.h}) implements forward
and backward passes:

\begin{lstlisting}[language=C++, caption={AutoDiffCircuit interface (autodiff.h)}]
// Forward: execute circuit with given params, return expectation
double forward(const std::vector<double>& params, const Observable& obs);

// Backward: parameter-shift gradient for all parameters
std::vector<double> backward(
    const std::vector<double>& params, const Observable& obs);
\end{lstlisting}

Supported observables: single-qubit Pauli $X,Y,Z$, and multi-qubit Pauli strings via
the \texttt{Observable} struct.

\subsubsection*{Distributed autodiff}

For circuits with many parameters, gradient computation parallelises trivially:
each parameter's two shifted evaluations are independent and can be dispatched to
separate MPI ranks. The implementation is in \texttt{src/distributed\_autodiff.cpp}.

\subsubsection*{JAX and PyTorch interfaces}

\begin{itemize}
  \item \textbf{JAX}~\cite{jax}: exposed via \texttt{python/qlret/jax\_interface.py};
        custom VJP registered with \texttt{jax.custom\_vjp} so \texttt{jax.grad},
        \texttt{jax.jit}, and \texttt{jax.vmap} all compose correctly.
  \item \textbf{PyTorch}~\cite{paszke2019pytorch}: a
        \texttt{torch.autograd.Function} subclass in
        \texttt{python/qlret/pytorch\_interface.py}, enabling LRET layers inside any
        PyTorch module with \texttt{.backward()}.
\end{itemize}

\subsection{PennyLane Device Plugin}

The \texttt{QLRETDevice} class (\texttt{python/qlret/pennylane\_device.py})
registers LRET as a native PennyLane device~\cite{bergholm2018pennylane} under the
name \texttt{"proxima.lret"}. Basic operation mapping (14 gates), Pauli and Hermitian
observables, tensor-product observables, parameter-shift gradients (via the module
above), finite-shot sampling, and probability measurement are implemented. A
differentiable-noise path required by gradient-based VQE/QAOA through noise channels
is ongoing work.

\subsection{Qiskit Provider}

The Qiskit~\cite{javadi2024qiskit} integration (\texttt{python/lret\_qiskit/}) exposes
LRET as a Qiskit provider with three pre-configured backends differing in truncation
tightness (\texttt{lret\_simulator}, \texttt{lret\_simulator\_accurate},
\texttt{lret\_simulator\_fast}). A circuit translator converts
\texttt{QuantumCircuit} objects to LRET JSON and a result converter maps LRET outputs
back to Qiskit \texttt{Result} objects. Provider registration follows the Qiskit
\texttt{BackendV2} interface. End-to-end execution (including IBM noise-model import)
passes its integration tests; broader hardware-representative noise validation is
future work.

\subsection{Cirq Integration}

The Cirq~\cite{cirq} integration (\texttt{python/lret\_cirq/}) provides
\texttt{LRETSimulator}, which inherits from \texttt{cirq.SimulatesSamples}. A
\texttt{CircuitTranslator} converts Cirq circuits to LRET JSON; a
\texttt{ResultConverter} maps LRET probability outputs back to Cirq \texttt{Result}
objects, preserving qubit maps and measurement keys. Parameterised circuits are
supported via Cirq's parameter resolver. Full native gate coverage and Cirq
noise-model import are ongoing work.

\subsection{Other Future Work}

\begin{itemize}
  \item Multi-node MPI performance benchmarking (\S\ref{sec:mpi}).
  \item GPU and multi-GPU performance benchmarking (\S\ref{sec:gpu}).
  \item Broader circuit-family coverage (variable depth, higher noise, structured
        circuits) for the LRET vs Cirq regime map.
  \item Continued optimisation of the row-parallel execution mode.
\end{itemize}

% ============================================================
\section{Proxima: AI-Driven Quantum Simulation Orchestration}
\label{sec:proxima}
% ============================================================

This section presents \textbf{Proxima}, the Python AI-orchestration layer that sits
above the LRET simulator and manages the full simulation lifecycle through a
conversational agent interface. Proxima is \emph{not} itself a quantum simulator; it
is a control-plane framework that abstracts the operational complexity of
LRET-based workflows (and five other quantum backends) behind natural-language
interaction, automated multi-step planning, and terminal-driven execution. The
treatment below mirrors the theoretical density of the preceding LRET sections: we
give the formal framing first (agentic control as a partially observable decision
process, hierarchical task decomposition, bandit-style model routing, circuit
fingerprinting), then describe how these theoretical objects are realised in code.
The subsystem is at a \emph{pre-release} stage of development; partial and planned
components are explicitly flagged throughout.

\subsection{Overview and the LRET Orchestration Problem}

Classical LRET simulation workflows require navigating multiple repositories,
managing Python virtual environments and C++ build systems, configuring framework
integrations (PennyLane~\cite{bergholm2018pennylane}, Cirq~\cite{cirq},
Qiskit~\cite{javadi2024qiskit}), and interpreting raw numerical output
(statevectors, density matrices, eigenvalue spectra, fidelity curves). The overhead
is significant, especially across the several LRET backend variants maintained on
different branches. Proxima was built to collapse this operational stack into a
single conversational interface.

The orchestration problem for quantum simulators decomposes into four structural
dimensions, which jointly motivate every design decision that follows:

\begin{enumerate}[leftmargin=1.5cm]
  \item \textbf{Installation complexity.} LRET and HPC backends
        (QuEST~\cite{jones2019quest}, qsim~\cite{qsim}) require C++17 compilation,
        CMake configuration, and optional GPU-SDK setup
        (cuQuantum~\cite{bayraktar2023cuquantum}, NVIDIA CUDA
        toolkits~\cite{nickolls2008cuda}). No unified installer covers all of
        them.
  \item \textbf{Workflow fragmentation.} A realistic campaign spans clone $\to$
        branch $\to$ configure $\to$ build $\to$ run $\to$ collect $\to$
        analyse $\to$ compare $\to$ export. Today these steps live in separate
        tools with no shared state, no recovery mechanism, and no end-to-end
        audit trail.
  \item \textbf{Result interpretation.} Raw simulator output demands domain
        expertise. An LLM-driven analysis layer bridges the gap between numerical
        artefacts (eigenvalue spectra, truncation ranks, Kraus-operator sets) and
        human-readable conclusions.
  \item \textbf{Backend heterogeneity.} Six frameworks expose six incompatible
        circuit formats, six gate-set conventions, and six result schemas;
        cross-backend comparison requires systematic normalisation.
\end{enumerate}

Proxima's intended role over the LRET simulation lifecycle is summarised in
Table~\ref{tab:proxima-lifecycle}. The second column is the outcome, not the
current state --- gaps are itemised in \S\ref{sec:proxima-gaps}.

\begin{table}[!h]
\centering
\small
\renewcommand{\arraystretch}{1.2}
\caption{Intended Proxima role across the LRET lifecycle.}
\label{tab:proxima-lifecycle}
\begin{tabular}{@{}ll@{}}
\toprule
\textbf{Phase} & \textbf{Proxima Role} \\
\midrule
Backend Acquisition & Clone LRET repo; checkout relevant branch. \\
Environment Setup   & Install Python/C++ dependencies; configure build. \\
Compilation         & Build native C++ extensions via CMake/setuptools. \\
Configuration       & Register the compiled backend with Proxima. \\
Execution           & Dispatch target simulation scripts or benchmarks. \\
Analysis            & Parse results; compare against FDM baselines; summarise. \\
Reporting           & Stream output during execution; surface final analysis in UI. \\
\bottomrule
\end{tabular}
\end{table}

\subsection{Theoretical Foundations of Agentic Orchestration}
\label{sec:proxima-theory}

Proxima's agent layer is a synthesis of three independent lineages in the
tool-using-LLM literature together with a control-theoretic framing of iterative
planning. The subsection makes the underlying assumptions explicit so that the
subsequent concrete design choices read as principled engineering rather than
ad-hoc convention.

\subsubsection*{The agent loop as a partially observable decision process}

Model a single agent turn as the tuple $(s_t, a_t, o_t, r_t)$ where $s_t$ is the
latent world state (filesystem, backend registry, session history, process-tree
state), $a_t$ is the tool call the model emits, $o_t$ is the observed tool result,
and $r_t$ is an implicit reward --- task completion as signalled by the user or by
a terminal tool such as \texttt{return\_result}. The model never observes $s_t$
directly, so the agent operates on a belief $b_t$ computed from (a) the
conversation prefix, (b) the running window of tool calls and results, and (c)
session grounding artefacts (current working directory, configured backend set,
recent checkpoints). Under this framing, the agent loop is the iteration
\[
  b_t \;\to\; a_t \;\to\; o_t \;\to\; b_{t+1}
\]
and the dual-model router (\S\ref{sec:proxima-dual-model}) is the policy-class
selector choosing between a low-cost belief-to-action mapping and a high-cost one.

\subsubsection*{Three lineages synthesised}

Proxima combines patterns from three families of prior work:

\begin{enumerate}[leftmargin=1.5cm]
  \item \textbf{ReAct}~\cite{yao2023react} interleaves an explicit \emph{Thought}
        step between observations and actions so the model reasons aloud about
        what tool to invoke and why. Proxima adopts the Observe--Think--Act cycle
        verbatim and streams \emph{Thought} tokens to the TUI for user
        observability.
  \item \textbf{Toolformer}~\cite{schick2023toolformer} casts tool invocation as
        structured generation against a typed schema rather than free-form text
        parsing. Proxima uses Pydantic-typed tool schemas and validates every
        emitted tool call against them before execution, which eliminates a large
        class of syntactically valid but semantically wrong tool calls.
  \item \textbf{Model Context Protocol (MCP)}~\cite{mcp2024} standardises the
        transport and discovery of tools across agent hosts. Proxima's tool
        registry is MCP-compatible by construction; the native MCP transport
        layer is deferred to the P2 roadmap (\S\ref{sec:proxima-roadmap}).
\end{enumerate}

\subsubsection*{Hierarchical task decomposition}

Context windows are finite. For long quantum workflows (clone $\to$ build $\to$
run $\to$ analyse) a naive strategy --- keep the full history in the main agent's
window --- exhausts model capacity within tens of tool calls and degrades
reasoning. The \emph{El Agente} quantum-chemistry agent of Zou \textit{et
al.}~\cite{zou2025elagente} addresses this with a hierarchical agent tree: a
strategic agent delegates sub-problems to specialised tactical agents which return
only \emph{summarised} feedback. This Markov-style context reduction keeps each
context window below roughly $10\%$ of the model's hard limit even across hundreds
of tool calls. Proxima's sub-agent architecture (\S\ref{sec:proxima-subagents}) is
a direct adaptation of this pattern, restricted in the current implementation to
read-only information-gathering roles. El Agente reports a success rate above
$87\%$ on university-level quantum-chemistry exercises without task-specific
hardcoding~\cite{zou2025elagente}; that empirical evidence that a hierarchical
LLM-driven agent can outperform rigid workflow pipelines on a scientifically rich
domain is the main justification for Proxima's LLM-mediated routing over a
hand-written dispatch table.

\subsubsection*{Multi-provider routing as a cost--quality bandit}

Routing a request to a small versus large model is a one-shot decision over an
expected-reward trade-off: the small tier is approximately an order of magnitude
cheaper and faster at real but workload-dependent accuracy cost. Proxima defers
the full contextual-bandit treatment; the current router uses a simple
complexity-score threshold --- a function of intent type, entity count,
conversation depth, and remaining token budget --- with a fixed cut point $c^{*}$.
Beyond $c^{*}$ the request escalates to the large tier.

\subsubsection*{Backend selection as circuit fingerprinting}

A quantum circuit can be characterised by a short feature vector
$\mathbf{f}(C) = (N, d, \mathbf{g}, \eta, \mathcal{O})$: qubit count $N$, depth
$d$, gate-type histogram $\mathbf{g}$, noise channel $\eta$, and required
observable $\mathcal{O}$. Proxima's backend selector
(\S\ref{sec:proxima-intelligence}) maps this fingerprint to a scored ranking over
registered adapters via deterministic feasibility filters (max-qubit, supported
gates, noise-model compatibility) followed by a soft score combining hardware
availability, memory feasibility, and observed historical latency. Where
cuQuantum GPU acceleration is feasible it typically
dominates~\cite{bayraktar2023cuquantum}; where LRET's low-rank structure applies
($p \lesssim 10^{-2}$, moderate entanglement) LRET~\cite{chen2021lret} dominates
on commodity hardware, as empirically confirmed by our §\ref{sec:empirical}
benchmark.

Together the Observe--Think--Act loop, hierarchical decomposition, bandit-style
routing, and circuit fingerprinting constitute the theoretical backbone on which
the concrete architectural choices below rest.

\subsection{Design Principles and Learning Sources}

\subsubsection*{Five design principles}

The architecture is governed by five principles, each traceable to a specific
failure mode observed during development:

\begin{enumerate}[leftmargin=1.5cm]
  \item \textbf{Controllability.} Every pipeline stage must support pause, resume,
        and rollback. No destructive operation proceeds without explicit user
        consent --- no exceptions.
  \item \textbf{Observability.} Execution state, resource usage, and LLM
        reasoning are surfaced in real time through the TUI and structured logs.
        ``Silent'' steps are the single largest source of perceived failure (see
        §\ref{sec:proxima-gaps}).
  \item \textbf{Composability.} Core components (pipeline, planner, executor,
        session manager) are independently instantiable and testable. No circular
        module dependencies.
  \item \textbf{Safety-by-default.} Scope-based tool permissions and consent
        gates prevent destructive operations from executing without explicit
        authorisation.
  \item \textbf{Backend neutrality.} The agent loop, pipeline, and UI are
        backend-agnostic. Backend-specific logic is confined to adapter modules
        under \texttt{src/proxima/backends/}.
\end{enumerate}

\subsubsection*{Learning sources}

Because Proxima sits at the intersection of two active research areas --- agentic
LLMs and numerical quantum simulation --- its design draws on two external
bodies of work rather than a single canonical reference.

\paragraph{Agentic architecture from \emph{El Agente}~\cite{zou2025elagente}.}
Four ideas were studied and adapted: (i) a cognitive-memory framework separating
global context, per-agent conversation history, grounding state, and a long-term
semantic/procedural store; (ii) hierarchical task decomposition with Markov-style
context reduction; (iii) a closed-loop error-recovery pattern alternating a
\emph{generation} phase (produce input files or commands) with a
\emph{verification} phase (parse errors and regenerate); (iv) transparent
action-trace logging so every tool invocation can be inspected and replayed
offline.

\paragraph{Engineering patterns from OpenCode AI~\cite{opencode2025}.} OpenCode
is a production-grade TypeScript code-agent under MIT licence
(\url{https://github.com/anomalyco/opencode}). Its modular monorepo layout
cleanly exposes tool registration, provider abstraction, scope-based permission
scopes, snapshot-based session recovery, and an MCP bridge. Proxima
reimplements analogous subsystems in Python, adopting OpenCode's
\emph{separation of provider abstraction from tool registry} as a first-class
architectural decision. Pydantic~2 replaces Zod for schema validation, SQLite
replaces the file-based session store, and the single-process async Python
design collapses OpenCode's client/server split because the Textual TUI already
provides an out-of-process UI surface.

\paragraph{Technical decisions.} The principal design choices and their
motivations:
\begin{itemize}[leftmargin=1.5cm]
  \item \emph{LLM-driven tool routing over a static dispatch table.} A static
        intent-to-tool table scales poorly once the tool count exceeds a few
        dozen; El Agente's high success rate on unseen tasks
        \emph{without} domain-specific hardcoding provides the empirical
        justification for routing through an LLM with a typed tool schema.
  \item \emph{Dual-model routing over single-model inference.} Separating
        low-complexity operations (entity extraction, parameter validation) from
        high-complexity reasoning (multi-step planning) allows the small-model
        tier to absorb most invocations.
  \item \emph{Permission-scoped sub-agents.} Backend compilation, package
        installation, and repository operations are dangerous when misdriven;
        OpenCode-style scope permissions restricted to a read-only subset for
        sub-agents bound the blast radius of a misbehaving model.
  \item \emph{Snapshot-based session recovery.} Long-running quantum workflows
        can fail at any step; El Agente's demonstration of sustained multi-turn
        operation via hierarchical memory motivates SQLite-backed sessions with
        stage-boundary checkpoints.
  \item \emph{LRET as the primary simulation target.} LRET's asymptotic advantage
        ($\mathcal{O}(V^2\cdot 2^N)$ vs.\ $\mathcal{O}(4^N)$) materialises only
        when a workflow selects the right backend per circuit instance, which
        motivates an intelligent dispatcher rather than a monolithic wrapper.
\end{itemize}

\subsection{Module Map and Core Data Structures}

\subsubsection*{Module map}

The implementation occupies $\sim\!370$ Python source files
($\approx 46{,}000$ lines of code under \texttt{src/proxima/**/*.py}). Its
top-level layout:

\begin{mdframed}[linecolor=gray!40,linewidth=0.5pt,
                 innertopmargin=6pt,innerbottommargin=6pt]
\dirtree{%
.1 src/proxima/.
.2 agent/ \DTcomment{tools, terminals, git, safety, planning}.
.3 dynamic\_tools/ \DTcomment{agent loop, NL processing, LLM bridge}.
.2 backends/ \DTcomment{six adapter files + LRET submodule}.
.3 lret/ \DTcomment{variants, installer, visualisation}.
.2 cli/ \DTcomment{Typer CLI, interactive REPL, formatters}.
.2 core/ \DTcomment{pipeline, planner, executor, state machine, session}.
.2 intelligence/ \DTcomment{LLM router, insights, backend selector}.
.2 resources/ \DTcomment{CPU/GPU/memory monitor, execution control}.
.2 tui/ \DTcomment{Textual app, screens, dialogs, wizards}.
.3 screens/ \DTcomment{main views}.
.3 dialogs/ \DTcomment{modal dialogs including backend wizard}.
.2 api/ \DTcomment{FastAPI app, middleware, routes, services}.
}
\end{mdframed}

User-facing surfaces are a Textual TUI~\cite{textual}, a Typer CLI~\cite{typer},
a FastAPI REST API~\cite{fastapi}, and an interactive REPL.

\subsubsection*{Core data structures}

\begin{description}[leftmargin=2.2cm,labelwidth=2.0cm,font=\ttfamily]
  \item[ExecutionState] Enum (\texttt{core/state.py}) covering the formal
        lifecycle \texttt{IDLE}, \texttt{PLANNING}, \texttt{READY},
        \texttt{RUNNING}, \texttt{PAUSED}, \texttt{COMPLETED}, \texttt{ABORTED},
        \texttt{ERROR}, \texttt{RECOVERING}. Illegal transitions raise
        \texttt{InvalidTransitionError}.
  \item[PersistedState] Serialisable snapshot of execution context: task DAG,
        partial results, resource handles, timestamps. Written at each
        checkpoint; restored on recovery.
  \item[ResourceHandle] RAII-style handle tracking acquired system resources
        (open processes, temporary files, GPU allocations). Released atomically
        on abort or rollback.
  \item[Session] Top-level container holding \texttt{session\_id},
        \texttt{execution\_history}, \texttt{checkpoints}, \texttt{context} dict,
        and optional \texttt{parent\_session\_id} for sub-agent sessions.
  \item[TaskNode] DAG vertex: \texttt{task\_id}, \texttt{task\_type},
        \texttt{parameters}, \texttt{dependencies}, \texttt{status},
        \texttt{result}.
  \item[BackendCapabilities] Descriptor: \texttt{max\_qubits},
        \texttt{supported\_gates}, \texttt{supported\_noise\_models},
        \texttt{gpu\_required}, \texttt{max\_shots}.
\end{description}

\subsection{Eight-Stage Execution Pipeline}

The pipeline (\texttt{core/pipeline.py}) implements a directed, fully
controllable flow:
\[
  \texttt{Parse}\,\to\,\texttt{Plan}\,\to\,\texttt{ResourceCheck}\,\to\,
  \texttt{Consent}\,\to\,\texttt{Execute}\,\to\,\texttt{Collect}\,\to\,
  \texttt{Analyse}\,\to\,\texttt{Export}.
\]

\begin{table}[!h]
\centering
\footnotesize
\renewcommand{\arraystretch}{1.2}
\caption{Eight-stage pipeline with control points.}
\label{tab:proxima-pipeline}
\begin{tabular}{@{}lp{4.0cm}p{5.3cm}l@{}}
\toprule
\textbf{Stage} & \textbf{Responsibility} & \textbf{Detail} & \textbf{Control} \\
\midrule
Parse         & Tokenise \& validate input       & Structural validation          & Pause \\
Plan          & Decompose into task DAG          & LLM-assisted decomposition     & Pause, Rollback \\
ResourceCheck & Verify CPU, RAM, GPU             & Pre-flight system checks       & Pause \\
Consent       & Summarise destructive ops        & Await explicit approval        & Gate \\
Execute       & Walk DAG; dispatch tools         & Per-task status tracking       & Pause, Resume, Rollback \\
Collect       & Aggregate results                & Schema normalisation           & Pause \\
Analyse       & Forward to LLM insight layer     & Structured summarisation       & Pause \\
Export        & Serialise JSON/CSV/Excel         & Session history update         & Pause \\
\bottomrule
\end{tabular}
\end{table}

Each stage emits events to \texttt{core/event\_bus.py}, consumed by the TUI's
reactive system for real-time progress display. Checkpoints are written at every
stage boundary, enabling resumption across process restarts.

\subsection{State Machine and Lifecycle Management}

\texttt{ExecutionStateMachine} (\texttt{core/state.py}) governs the lifecycle via
a transition table with entry and exit callbacks per state. The ten validated
transitions are:

\begin{center}
\small
\renewcommand{\arraystretch}{1.15}
\begin{tabular}{@{}lll@{}}
\toprule
\textbf{From} & \textbf{To} & \textbf{Trigger} \\
\midrule
\texttt{IDLE}       & \texttt{PLANNING}  & user command \\
\texttt{PLANNING}   & \texttt{READY}     & DAG finalised, resources confirmed \\
\texttt{READY}      & \texttt{RUNNING}   & user consent obtained \\
\texttt{RUNNING}    & \texttt{PAUSED}    & user interrupt / resource limit \\
\texttt{PAUSED}     & \texttt{RUNNING}   & user resume \\
\texttt{RUNNING}    & \texttt{ERROR}     & unrecoverable tool failure \\
\texttt{ERROR}      & \texttt{RECOVERING} & recovery strategy identified \\
\texttt{RECOVERING} & \texttt{RUNNING}   & recovery succeeded \\
\texttt{RUNNING}    & \texttt{COMPLETED} & all tasks complete \\
\texttt{RUNNING}    & \texttt{ABORTED}   & user cancellation or rollback \\
\bottomrule
\end{tabular}
\end{center}

State is persisted to disk at every transition boundary, enabling recovery across
process restarts.

\subsection{DAG Task Planning, Execution, and Recovery}

\textbf{Planner} (\texttt{core/planner.py}) receives a natural-language command or
structured circuit request and produces a DAG of \texttt{TaskNode} objects.
LLM-assisted decomposition handles ambiguous multi-step commands; each node
declares explicit dependencies, and parallel-safe tasks are identified
automatically.

\textbf{Executor} (\texttt{core/executor.py}) topologically sorts the DAG and
walks it under the \texttt{ExecutionStateMachine}. Per-task status is tracked
independently, enabling partial re-execution: a failed subtask can be retried
without re-running completed predecessors. Resource allocation and cleanup are
managed by \texttt{ResourceHandle} objects attached to each \texttt{TaskNode}.

\textbf{Error recovery.} On task failure, \texttt{AgentErrorHandler} classifies
the error (network, compiler, permission, timeout, unknown), selects a recovery
strategy (retry with backoff, suggest manual intervention, rollback to
checkpoint), and transitions the state machine accordingly.

\subsection{Session Management and Checkpointing}

Each user interaction creates or resumes a \texttt{Session} object stored in
SQLite. Sessions persist conversation history, execution history, and checkpoints
across process restarts. Multi-session concurrency is supported; sessions are
isolated by \texttt{session\_id}.

A \texttt{Checkpoint} captures the full \texttt{PersistedState} at a given
execution point. Checkpoints are created automatically at pipeline-stage
boundaries and on demand. Rollback restores the most recent checkpoint and
reverses any filesystem changes tracked by \texttt{FileResource} handles.

When the accumulated token count approaches the LLM context limit,
\texttt{AgentSessionManager} invokes auto-summarisation: a summary LLM call
condenses prior turns into a compact context block, transparently continuing the
session. Sub-agents launched during a parent session create child sessions with
\texttt{parent\_session\_id} set, forming hierarchical session trees.

\subsection{ReAct-Style Agent Loop}
\label{sec:proxima-loop}

The agent loop (\texttt{agent/dynamic\_tools/agent\_loop.py}) implements the
ReAct~\cite{yao2023react} observe--think--act cycle as an async coroutine with
transparent action-trace logging after El Agente~\cite{zou2025elagente}:

\begin{lstlisting}[language=Python, caption={Agent loop pseudocode}]
observe(user_input)
  |-- NL parsing      --> Intent object
  |-- Session context extraction
  `-- High-confidence intent --> direct dispatch (skip Think)

think(intent)
  |-- LLM routing     --> model tier selection
  |-- Tool selection  --> ToolCall specification
  `-- Chain-of-thought streamed to TUI

act(tool_call)
  |-- ToolPermissionManager.check_permission()
  |-- ToolRegistry.execute()      --> streaming output
  `-- ResultAnalyzer.analyze()    --> StructuredResult

observe(result) --> repeat or return
\end{lstlisting}

\textbf{Direct dispatch.} Intents classified above a confidence threshold bypass
LLM reasoning and are dispatched directly to the corresponding tool, reducing
median latency for common operations (file reads, git status, backend listing)
to sub-100\,ms.

\textbf{Streaming.} LLM chain-of-thought tokens and tool output are streamed
incrementally to the TUI AI Assistant screen via the event bus, providing
real-time feedback without blocking the Textual event loop.

\textbf{Context management.} The agent loop tracks accumulated token count
across turns; when the budget is $\sim 80\%$ consumed,
\texttt{AgentSessionManager.summarize()} is invoked transparently before the
next LLM call.

\subsection{Tool Permission and Safety}

\texttt{ToolPermissionManager} (\texttt{agent/tool\_permissions.py}) implements
scope-based access control. Each tool is tagged with a \texttt{PermissionScope};
the manager enforces the active session's permission set before every tool
invocation. The hierarchy satisfies
\[
  \texttt{READ\_ONLY}\;\subset\;\texttt{WRITE\_LOCAL}\;\subset\;
  \texttt{WRITE\_REMOTE}\;\subset\;\texttt{DESTRUCTIVE}\;\subset\;\texttt{ADMIN}.
\]

\textbf{\texttt{SafetyManager}} detects destructive operations (file deletion,
branch force-push, database drop, process kill) and routes them to the consent
gate before execution, maintaining a rollback stack of reversible operations.
\textbf{\texttt{ConsentManager}} presents structured consent dialogs --- operation
summary, affected resources, reversibility assessment, explicit approve/deny
buttons --- with every decision logged to the audit trail.
\textbf{\texttt{AdminPrivilegeHandler}} manages elevation requests for operations
requiring system-level permissions (Windows UAC, Unix sudo).

\subsection{Sub-Agent Architecture}
\label{sec:proxima-subagents}

Sub-agents (\texttt{agent/sub\_agent.py}) are lightweight read-only agent
instances launched by the main loop to parallelise information gathering without
exposing write tools. The pattern --- strategic reasoning at the parent level
with summarised results returned from tactical children --- follows El
Agente~\cite{zou2025elagente} and the sub-agent conventions of
OpenCode~\cite{opencode2025}.

\texttt{SubAgentConfig} declares: \texttt{name}; \texttt{allowed\_tools}
(read-only subset: \texttt{read\_file}, \texttt{search\_files}, \texttt{git\_log},
\texttt{list\_directory}, \texttt{grep\_content}); \texttt{model\_preference}
(``small''); \texttt{max\_iterations}; and \texttt{auto\_approve\_permissions}.
\texttt{SubAgent.run(prompt)} executes a simplified single-goal loop and returns
a \texttt{SubAgentResult} to the parent. Each run creates a child session with
\texttt{parent\_session\_id} set, supporting hierarchical session trees.

\subsection{Dual-Model Router}
\label{sec:proxima-dual-model}

\texttt{DualModelRouter} assigns LLM roles to two tiers:

\begin{center}
\small
\renewcommand{\arraystretch}{1.2}
\begin{tabular}{@{}lp{9cm}@{}}
\toprule
\textbf{Tier} & \textbf{Responsibilities} \\
\midrule
Large (Claude Opus~\cite{anthropic_claude}, GPT-4~\cite{openai_gpt4}) &
  complex multi-step reasoning; plan decomposition; ambiguous intent
  classification; domain-aware result analysis. \\
Small (Claude Haiku, GPT-3.5/4o-mini, Ollama~\cite{ollama} local models) &
  entity extraction; parameter validation; progress narration; sub-agent tasks. \\
\bottomrule
\end{tabular}
\end{center}

Routing is per-invocation based on a complexity score derived from intent type,
entity count, conversation depth, and remaining token budget. If the preferred
provider is unavailable, the router falls back transparently.

\subsection{Intelligence Layer}
\label{sec:proxima-intelligence}

\subsubsection*{Multi-provider LLM routing}

\texttt{LLMRouter} (\texttt{intelligence/llm\_router.py}) abstracts three
providers --- Anthropic Claude~\cite{anthropic_claude}, OpenAI
GPT-4~\cite{openai_gpt4}, and the local Ollama runtime~\cite{ollama} --- behind a
unified async interface. Provider-fallback covers rate limits and quota
exhaustion; token-budget management triggers summarisation; all providers are
accessed via native streaming APIs with tokens forwarded to the TUI event bus.
A dual-model strategy routes high-complexity planning to the large tier and
high-frequency low-complexity operations (entity extraction, parameter
validation) to the small tier; internal measurements indicate $60\text{--}70\%$
token-consumption reduction on representative simulation-campaign traces with no
measurable quality regression, though these numbers are indicative and a
reproducible harness is part of the near-term roadmap.

\subsubsection*{Backend auto-selection}

\texttt{BackendSelector} (\texttt{intelligence/selector.py}) scores registered
backends against a six-factor evaluation --- the concrete realisation of the
circuit fingerprint $\mathbf{f}(C)$ introduced in \S\ref{sec:proxima-theory}:

\begin{enumerate}[leftmargin=1.5cm]
  \item Qubit-count feasibility: reject backends with
        $\texttt{max\_qubits} < \texttt{circuit\_qubits}$.
  \item Gate compatibility: penalise or eliminate circuits containing
        unsupported gates.
  \item Noise-model support: match the requested noise model against declared
        backend capabilities.
  \item Hardware availability: score GPU-required backends negatively when no
        CUDA-capable device is detected.
  \item Memory estimation: compare estimated memory against available RAM.
  \item Historical performance: retrieve past execution times for similar
        circuits from the SQLite telemetry store.
\end{enumerate}

The top-ranked backend is selected by default; the user may override via explicit
parameter or TUI selection.

\subsubsection*{Result analysis and insight generation}

\texttt{ResultAnalyzer} (\texttt{agent/dynamic\_tools/result\_analyzer.py})
extracts structured metrics from raw backend output: fidelity (trace distance or
state fidelity against reference), wall-clock and CPU time per gate/layer/total,
peak RAM and GPU VRAM, LRET truncation rank $V$, and noise sensitivity $F(p)$.
\texttt{InsightGenerator} (\texttt{intelligence/insights.py}) feeds the structured
metrics to the LLM with a domain-specific system prompt, producing human-readable
summaries: trend identification, anomaly flagging, cross-backend comparisons,
and parameter-tuning recommendations.

\subsection{Multi-Backend Adapter Architecture}

All backend adapters extend \texttt{BaseBackendAdapter} and implement a standard
interface:

\begin{lstlisting}[language=Python, caption={\texttt{BaseBackendAdapter} interface}]
class BaseBackendAdapter:
    def execute(self, circuit, shots, noise_model=None) -> BackendResult
    def get_capabilities(self) -> BackendCapabilities
    def estimate_performance(self, circuit) -> PerformanceEstimate
    def validate_circuit(self, circuit) -> ValidationResult
    def health_check(self) -> HealthStatus
\end{lstlisting}

\texttt{BackendRegistry} maintains the live registry, performs periodic health
checks, and provides the comparison matrix used by the selector. Circuit-format
conversion (\texttt{backends/conversion.py}) and result normalisation
(\texttt{backends/normalization.py}) are applied transparently; GPU memory
allocation is managed centrally by \texttt{GPUMemoryManager}.

\begin{table}[!h]
\centering
\small
\renewcommand{\arraystretch}{1.2}
\caption{Registered backends and their adapter status.}
\label{tab:proxima-backends}
\begin{tabular}{@{}lcccl@{}}
\toprule
\textbf{Backend} & \textbf{Max qubits} & \textbf{GPU} & \textbf{Status} & \textbf{Notes} \\
\midrule
LRET                                      & variable & no       & primary target & custom low-rank simulator \\
Cirq~\cite{cirq}                          & 30+      & no       & functional     & Google ecosystem \\
Qiskit Aer~\cite{javadi2024qiskit}        & 30+      & optional & functional     & Windows build caveats \\
cuQuantum~\cite{bayraktar2023cuquantum}   & 30+      & required & functional     & requires NVIDIA SDK \\
qsim~\cite{qsim}                          & 30+      & optional & partial        & build not automated \\
QuEST~\cite{jones2019quest}               & 30+      & optional & partial        & multi-node HPC capable \\
\bottomrule
\end{tabular}
\end{table}

A PennyLane~\cite{bergholm2018pennylane} device plugin \texttt{proxima.lret}
exposes the simulator as a native \texttt{qml.device(...)} target; gate mapping
and non-differentiable circuit execution are functional, while differentiable
noise simulation required for gradient-based VQE/QAOA/QNN workflows is part of
the roadmap.

\subsection{Orchestration-Overhead Disclaimer and Benchmarks}

\begin{infobox}
\textbf{Disclaimer.} Orchestration benchmark figures reported in earlier
Proxima internals (e.g., ``$2.5\text{--}5\times$ workflow speedup over manual
execution'', ``$15\text{--}40$\,ms per-execution overhead'', ``$60\text{--}70\%$
token savings via dual-model routing'') are \emph{illustrative targets pending
formal benchmark runs}, not results from a single end-to-end reproducible
harness. They are quoted in this section only to communicate design intent and
indicative measured behaviour.
\end{infobox}

For single-circuit single-backend execution with no post-processing, direct
backend invocation is faster than Proxima by the orchestration-overhead margin.
Proxima's value is in \emph{multi-step} workflows (clone$\to$configure$\to$
build$\to$run$\to$analyse), \emph{cross-backend} comparison, and \emph{novel
practitioner} onboarding --- the three settings where DAG parallelism, schema
normalisation, and LLM-driven insight generation amortise the control-plane
cost.

\subsection{Current Implementation Status and Gaps}
\label{sec:proxima-gaps}

Proxima is pre-release (v0.1-dev). The control-plane scaffolding (pipeline,
planner, state machine, session management, permission system, adapter layer) is
substantially in place, but several execution paths remain incomplete or only
partially wired. Honest enumeration:

\begin{itemize}[leftmargin=1.5cm]
  \item \textbf{C++ backend build pipeline.} \texttt{backend\_builder.py} defines
        the clone$\to$cmake$\to$compile$\to$register interface, but the full
        automation does not yet run end-to-end on Windows and is not connected to
        the TUI backend wizard.
  \item \textbf{Terminal-to-TUI output bridge.} \texttt{terminal\_tui\_bridge.py}
        exists as a bridge component but the live connection to the Execution
        screen does not reliably stream subprocess output in real time --- users
        perceive long tool calls as silent, eroding trust.
  \item \textbf{Post-execution auto-analysis.} Result Tab does not automatically
        populate with LLM-generated analysis after a run completes; the user must
        request it.
  \item \textbf{Complex multi-step NL execution.} \texttt{robust\_nl\_processor.py}
        and \texttt{intent\_classifier.py} produce plans, but chaining NL parsing
        $\to$ plan $\to$ execute $\to$ result display as a single coherent
        workflow is not reliable; mid-plan failures do not propagate cleanly to
        the user.
  \item \textbf{Monolithic source files.} \texttt{intent\_tool\_bridge.py}
        ($272$\,KB), \texttt{llm\_router.py} ($311$\,KB), and
        \texttt{agent\_ai\_assistant.py} ($381$\,KB) impair maintainability.
  \item \textbf{Type-safety suppression.} \texttt{mypy} is configured to ignore
        errors in \texttt{proxima.tui.*}, \texttt{proxima.backends.conversion},
        \texttt{proxima.backends.normalization}, \texttt{proxima.backends.lret.*},
        and \texttt{proxima.data.export}; suppressions mask real type errors.
  \item \textbf{Windows-specific issues.} Qiskit~Aer requires MSVC or MinGW and
        may silently fail without one; the Docker image is not published; some
        internal path constructions use hardcoded forward slashes.
\end{itemize}

\subsection{Roadmap}
\label{sec:proxima-roadmap}

Tasks are ordered by priority. \textcolor{p1color}{\textbf{P1 = blocking}},
\textcolor{p2color}{\textbf{P2 = high}}, \textcolor{p3color}{\textbf{P3 = medium}}.

\begin{table}[!h]
\centering
\footnotesize
\renewcommand{\arraystretch}{1.2}
\caption{\textcolor{p1color}{P1 --- critical blockers.}}
\begin{tabular}{@{}p{3.8cm}p{10.8cm}@{}}
\toprule
\textbf{Task} & \textbf{Suggested Approach} \\
\midrule
Wire terminal to Execution Tab           & connect \texttt{process\_streaming.py} events via \texttt{event\_bus.py} to a Textual \texttt{RichLog} widget \\
Implement C++ build pipeline             & compiler detection (CMake, MSVC, GCC), cmake configure/build, post-build tests, error surfacing \\
Auto-register compiled backend           & update \texttt{backends/registry.py} + config entry; success notification \\
Close NL\,$\to$\,plan\,$\to$\,execute gap & fully delegate plan execution to \texttt{plan\_executor.py} with per-step TUI callbacks \\
Post-execution auto-analysis              & on completion, invoke \texttt{result\_analyzer.py} + \texttt{insights.py} and push to Result Tab \\
\bottomrule
\end{tabular}
\end{table}

\begin{table}[!h]
\centering
\footnotesize
\renewcommand{\arraystretch}{1.2}
\caption{\textcolor{p2color}{P2 --- high priority.}}
\begin{tabular}{@{}p{3.8cm}p{10.8cm}@{}}
\toprule
\textbf{Task} & \textbf{Suggested Approach} \\
\midrule
LRET installer cancellation      & cancellation token in \texttt{installer\_dialog.py} propagated to the subprocess \\
Backend-config dialog pre-fill   & load existing backend config before opening dialog \\
C++ dependency management        & detect CMake, compiler toolchains, CUDA SDK; install instructions when missing \\
Authenticated GitHub operations  & HTTPS token / SSH key handling; credentials stored via \texttt{keyring} \\
PennyLane differentiable device  & finish differentiable noise simulation in \texttt{backends/lret/pennylane\_device.py} \\
\bottomrule
\end{tabular}
\end{table}

\begin{table}[!h]
\centering
\footnotesize
\renewcommand{\arraystretch}{1.2}
\caption{\textcolor{p3color}{P3 --- medium priority.}}
\begin{tabular}{@{}p{3.8cm}p{10.8cm}@{}}
\toprule
\textbf{Task} & \textbf{Suggested Approach} \\
\midrule
Decompose oversized files   & split \texttt{intent\_tool\_bridge.py}, \texttt{llm\_router.py}, \texttt{agent\_ai\_assistant.py} into cohesive sub-modules \\
Re-enable \texttt{mypy}     & fix type errors in suppressed modules progressively, starting with the smallest \\
Multi-terminal split-pane   & add tabbed/split-panel terminal view for concurrent process monitoring \\
Publish Docker image        & GitHub Actions workflow to build and push to GHCR on release tag \\
Native MCP transport        & port \texttt{packages/opencode/src/mcp/} analogue to \texttt{src/proxima/mcp/} \\
\bottomrule
\end{tabular}
\end{table}

\subsection{Key Learnings}

Six learnings crystallised during development that future contributors (and
readers of this report) should inherit:

\begin{enumerate}[leftmargin=1.5cm]
  \item \emph{LLM-driven routing only outperforms a static dispatch table when
        tool schemas are strict.} Loosely typed tool parameters let the LLM emit
        syntactically valid but semantically wrong calls that fail downstream
        and erode the agent's perceived reliability.
  \item \emph{Monolithic files impose a recurring tax.} A single 381-KB TUI
        screen costs more in LLM context-per-call than five smaller modules of
        equivalent total size --- oversized files impair human maintenance
        \emph{and} agent self-modification.
  \item \emph{Most observed ``failure cases'' reduce to observability gaps.}
        The TUI streaming bridge, post-execution auto-analysis, and mid-plan
        error surfacing are UI-observability problems more than algorithmic
        defects; users perceive an algorithm as broken when they cannot see it
        working.
  \item \emph{Reusing OpenCode's provider-abstraction pattern saved substantial
        work.} The clean separation of \emph{provider} from \emph{tool} in
        OpenCode~\cite{opencode2025} allowed Anthropic, OpenAI, and Ollama
        support to share one streaming and fallback implementation rather than
        three bespoke ones.
  \item \emph{El Agente's hierarchical memory is essential for long sessions.}
        Once a Proxima session exceeds roughly $50{,}000$ tokens, single-agent
        designs degrade noticeably; hierarchical sub-agents with summarised
        returns restore reliability~\cite{zou2025elagente}.
  \item \emph{Truthful documentation matters as much as correct code.}
        Documentation that over-claims capability erodes user trust on first
        encounter with an unmet expectation; this section therefore marks
        partial and planned work explicitly throughout, consistent with the
        tone set in the rest of the report.
\end{enumerate}

% ============================================================
\section{Conclusion}
\label{sec:conclusion}
% ============================================================

We have presented an engineering treatment of the Low-Rank Eigenvalue Truncation
algorithm of Chen, Farquhar and Parrish~\cite{chen2021lret}. The algorithm itself is
prior work; our contribution is a production C++ implementation with:
(i) 21 gate types and 11 physical noise channels, each cited to its primary source;
(ii) five OpenMP~\cite{dagummenon1998} execution modes with formally stated
algorithms and cost models in Foster's framework~\cite{foster1995}, supplemented by
gate fusion after qsim~\cite{qsim} and circuit stratification after
$\mathsf{t\vert ket\rangle}$~\cite{sivarajah2020tket};
(iii) MPI~\cite{gropp1999mpi} and GPU (cuQuantum~\cite{bayraktar2023cuquantum} +
custom CUDA~\cite{nickolls2008cuda}) code paths with smoke tests passing;
(iv) Python bindings and in-progress integrations with
PennyLane~\cite{bergholm2018pennylane}, Qiskit~\cite{javadi2024qiskit}, and
Cirq~\cite{cirq};
(v) a parameter-shift autodiff module~\cite{mitarai2018,schuld2019} with JAX and
PyTorch interfaces; and
(vi) the Proxima orchestration framework wrapping all of the above behind a
conversational agent.

Our empirical validation is deliberately scoped. In a head-to-head against Cirq's
density-matrix simulator~\cite{cirq} on random dense depth-13 circuits with
$p=10^{-3}$ depolarising noise, LRET reaches a $59.1\times$ speedup at $N=14$ with
state fidelity $\ge 0.986$, and continues to execute at $N=16, 18$ where the
full-density-matrix reference exceeds host memory. Final ranks stay at or below
$0.058\%$ of the Hilbert dimension at $N=18$, supporting the low-rank assumption
of~\cite{chen2021lret}. All other subsystems (MPI, GPU, QEC, autodiff) are verified
at the smoke-test level only; performance characterisation of these paths is
deferred to future work, as are broader circuit families, higher noise rates, and
multi-node scaling studies.

We believe this scoping is itself the contribution: a faithful, well-attributed,
reproducible implementation of a published algorithm, with one focused empirical
result and a clear roadmap for the remaining measurement work.

\section*{Acknowledgements}

The LRET algorithm we implement is due to Chen, Farquhar \& Parrish~\cite{chen2021lret}.
The gate-fusion pass is inspired by Google's qsim~\cite{qsim} (Apache 2.0).
The MPI row-distribution strategy is inspired by QuEST~\cite{jones2019quest} (MIT).
The Proxima agent architecture draws on El Agente~\cite{zou2025elagente} and
OpenCode~\cite{opencode2025} (MIT).

% ============================================================
\begin{thebibliography}{99}
% ============================================================

\bibitem{chen2021lret}
  Y.~Chen, C.~Farquhar, R.~M.~Parrish.
  \textit{Low-rank density-matrix evolution for noisy quantum circuits}.
  npj Quantum Information \textbf{7}, 61 (2021).

\bibitem{nielsenchuang2010}
  M.~A.~Nielsen, I.~L.~Chuang.
  \textit{Quantum Computation and Quantum Information}, 10th Anniversary Ed.
  Cambridge University Press (2010).

\bibitem{kraus1983}
  K.~Kraus.
  \textit{States, Effects, and Operations: Fundamental Notions of Quantum Theory}.
  Springer Lecture Notes in Physics \textbf{190} (1983).

\bibitem{choi1975}
  M.-D.~Choi.
  \textit{Completely positive linear maps on complex matrices}.
  Linear Algebra and its Applications \textbf{10}, 285--290 (1975).

\bibitem{breuerpetruccione2002}
  H.-P.~Breuer, F.~Petruccione.
  \textit{The Theory of Open Quantum Systems}. Oxford University Press (2002).

\bibitem{lindblad1976}
  G.~Lindblad.
  \textit{On the generators of quantum dynamical semigroups}.
  Commun.\ Math.\ Phys.\ \textbf{48}, 119--130 (1976).

\bibitem{krantz2019}
  P.~Krantz, M.~Kjaergaard, F.~Yan, T.~P.~Orlando, S.~Gustavsson, W.~D.~Oliver.
  \textit{A quantum engineer's guide to superconducting qubits}.
  Appl.\ Phys.\ Rev.\ \textbf{6}, 021318 (2019).

\bibitem{mcewen2021leakage}
  M.~McEwen \textit{et al.}
  \textit{Removing leakage-induced correlated errors in superconducting quantum
  error correction}.
  Nature Communications \textbf{12}, 1761 (2021).

\bibitem{miao2023leakage}
  K.~C.~Miao \textit{et al.}
  \textit{Overcoming leakage in quantum error correction}.
  Nature Physics \textbf{19}, 1780--1786 (2023).

\bibitem{gottesman1997}
  D.~Gottesman.
  \textit{Stabilizer Codes and Quantum Error Correction}.
  Ph.D.\ thesis, Caltech (1997). arXiv:quant-ph/9705052.

\bibitem{aaronson2004chpf}
  S.~Aaronson, D.~Gottesman.
  \textit{Improved simulation of stabilizer circuits}.
  Phys.\ Rev.\ A \textbf{70}, 052328 (2004).

\bibitem{zhang2017correlated}
  J.~Zhang \textit{et al.}
  \textit{Correlated quantum errors and their impact on quantum error correction}.
  Phys.\ Rev.\ Lett.\ \textbf{118}, 230502 (2017).

\bibitem{clader2021}
  B.~D.~Clader \textit{et al.}
  \textit{Quantum resources required to block-encode a matrix of classical data}.
  arXiv:2206.03505 (2022).

\bibitem{mitarai2018}
  K.~Mitarai, M.~Negoro, M.~Kitagawa, K.~Fujii.
  \textit{Quantum circuit learning}.
  Phys.\ Rev.\ A \textbf{98}, 032309 (2018).

\bibitem{schuld2019}
  M.~Schuld, V.~Bergholm, C.~Gogolin, J.~Izaac, N.~Killoran.
  \textit{Evaluating analytic gradients on quantum hardware}.
  Phys.\ Rev.\ A \textbf{99}, 032331 (2019).

\bibitem{banchi2021}
  L.~Banchi, G.~E.~Crooks.
  \textit{Measuring analytic gradients of general quantum evolution with the
  stochastic parameter shift rule}.
  Quantum \textbf{5}, 386 (2021).

\bibitem{jones2019quest}
  T.~Jones, A.~Brown, I.~Bush, S.~C.~Benjamin.
  \textit{QuEST and high performance simulation of quantum computers}.
  Scientific Reports \textbf{9}, 10736 (2019).

\bibitem{qsim}
  Google Quantum AI and collaborators.
  \textit{qsim: Schr\"odinger and Schr\"odinger--Feynman simulators for quantum
  circuits}.
  \url{https://github.com/quantumlib/qsim} (accessed 2026).

\bibitem{hanersteiger2017}
  T.~H\"aner, D.~S.~Steiger.
  \textit{0.5 petabyte simulation of a 45-qubit quantum circuit}.
  In \emph{Proceedings of SC'17} (2017).

\bibitem{khammassi2017qx}
  N.~Khammassi, I.~Ashraf, X.~Fu, C.~G.~Almudever, K.~Bertels.
  \textit{QX: a high-performance quantum computer simulation platform}.
  DATE (2017).

\bibitem{sivarajah2020tket}
  S.~Sivarajah, S.~Dilkes, A.~Cowtan, W.~Simmons, A.~Edgington, R.~Duncan.
  \textit{t$|$ket$\rangle$: a retargetable compiler for NISQ devices}.
  Quantum Science and Technology \textbf{6}, 014003 (2020).

\bibitem{cowtan2019}
  A.~Cowtan, S.~Dilkes, R.~Duncan, A.~Krajenbrink, W.~Simmons, S.~Sivarajah.
  \textit{On the qubit routing problem}. TQC 2019.

\bibitem{zulehner2018qmdd}
  A.~Zulehner, R.~Wille.
  \textit{Advanced simulation of quantum computations}.
  IEEE TCAD \textbf{38}, 848--859 (2018).

\bibitem{cormen2009}
  T.~H.~Cormen, C.~E.~Leiserson, R.~L.~Rivest, C.~Stein.
  \textit{Introduction to Algorithms}, 3rd ed. MIT Press (2009).

\bibitem{foster1995}
  I.~Foster.
  \textit{Designing and Building Parallel Programs}. Addison-Wesley (1995).

\bibitem{dagummenon1998}
  L.~Dagum, R.~Menon.
  \textit{OpenMP: an industry-standard API for shared-memory programming}.
  IEEE Computational Science and Engineering \textbf{5}, 46--55 (1998).

\bibitem{gropp1999mpi}
  W.~Gropp, E.~Lusk, A.~Skjellum.
  \textit{Using MPI: Portable Parallel Programming with the Message-Passing
  Interface}. MIT Press (1999).

\bibitem{pacheco1996pp}
  P.~Pacheco.
  \textit{Parallel Programming with MPI}. Morgan Kaufmann (1996).

\bibitem{valiant1990bsp}
  L.~G.~Valiant.
  \textit{A bridging model for parallel computation}.
  Comm.\ ACM \textbf{33}, 103--111 (1990).

\bibitem{amdahl1967}
  G.~M.~Amdahl.
  \textit{Validity of the single processor approach to achieving large scale
  computing capabilities}. AFIPS Conf.\ Proc.\ \textbf{30}, 483--485 (1967).

\bibitem{wilkinson2005}
  B.~Wilkinson, M.~Allen.
  \textit{Parallel Programming: Techniques and Applications Using Networked
  Workstations and Parallel Computers}, 2nd ed. Pearson (2005).

\bibitem{rabenseifner2009}
  R.~Rabenseifner, G.~Hager, G.~Jost.
  \textit{Hybrid MPI/OpenMP parallel programming on clusters of multi-core SMP
  nodes}. PDP 2009.

\bibitem{hennessypatterson2017}
  J.~L.~Hennessy, D.~A.~Patterson.
  \textit{Computer Architecture: A Quantitative Approach}, 6th ed. Morgan Kaufmann
  (2017).

\bibitem{goto2008anatomy}
  K.~Goto, R.~A.~van de Geijn.
  \textit{Anatomy of high-performance matrix multiplication}.
  ACM Transactions on Mathematical Software \textbf{34}, 12 (2008).

\bibitem{goluvanloan2013}
  G.~H.~Golub, C.~F.~Van Loan.
  \textit{Matrix Computations}, 4th ed. Johns Hopkins (2013).

\bibitem{eigen2010}
  G.~Guennebaud, B.~Jacob \textit{et al.}
  \textit{Eigen v3}. \url{https://eigen.tuxfamily.org} (2010).

\bibitem{pybind11}
  W.~Jakob, J.~Rhinelander, D.~Moldovan.
  \textit{pybind11 --- seamless operability between C++11 and Python}.
  \url{https://github.com/pybind/pybind11} (accessed 2026).

\bibitem{cmake}
  K.~Martin, B.~Hoffman.
  \textit{Mastering CMake}. Kitware (2007).

\bibitem{nickolls2008cuda}
  J.~Nickolls, I.~Buck, M.~Garland, K.~Skadron.
  \textit{Scalable parallel programming with CUDA}. ACM Queue \textbf{6}, 40--53 (2008).

\bibitem{lindholm2008tesla}
  E.~Lindholm, J.~Nickolls, S.~Oberman, J.~Montrym.
  \textit{NVIDIA Tesla: A unified graphics and computing architecture}.
  IEEE Micro \textbf{28}, 39--55 (2008).

\bibitem{volkovdemmel2008}
  V.~Volkov, J.~W.~Demmel.
  \textit{Benchmarking GPUs to tune dense linear algebra}. SC'08.

\bibitem{nvidiavolta2017}
  NVIDIA Corporation.
  \textit{NVIDIA Tesla V100 GPU Architecture} (Whitepaper, 2017).

\bibitem{bayraktar2023cuquantum}
  H.~Bayraktar \textit{et al.}
  \textit{cuQuantum SDK: A high-performance library for accelerating quantum
  science}. IEEE QCE 2023.

\bibitem{nccl}
  NVIDIA Collective Communications Library (NCCL).
  \url{https://github.com/NVIDIA/nccl} (accessed 2026).

\bibitem{bergholm2018pennylane}
  V.~Bergholm \textit{et al.}
  \textit{PennyLane: Automatic differentiation of hybrid quantum--classical
  computations}. arXiv:1811.04968 (2018).

\bibitem{javadi2024qiskit}
  A.~Javadi-Abhari \textit{et al.}
  \textit{Quantum computing with Qiskit}. arXiv:2405.08810 (2024).

\bibitem{cirq}
  Cirq Developers.
  \textit{Cirq: A Python framework for NISQ algorithms}.
  \url{https://github.com/quantumlib/Cirq} (accessed 2026).

\bibitem{jax}
  J.~Bradbury \textit{et al.}
  \textit{JAX: composable transformations of Python+NumPy programs}.
  \url{https://github.com/google/jax} (2018--).

\bibitem{paszke2019pytorch}
  A.~Paszke \textit{et al.}
  \textit{PyTorch: An imperative style, high-performance deep learning library}.
  NeurIPS (2019).

\bibitem{orus2014tn}
  R.~Or\'us.
  \textit{A practical introduction to tensor networks: Matrix product states and
  projected entangled pair states}.
  Annals of Physics \textbf{349}, 117--158 (2014).

\bibitem{schollwock2011dmrg}
  U.~Schollw\"ock.
  \textit{The density-matrix renormalization group in the age of matrix product
  states}.
  Annals of Physics \textbf{326}, 96--192 (2011).

\bibitem{zou2025elagente}
  Y.~Zou \textit{et al.}
  \textit{El Agente: An autonomous agent for quantum chemistry}.
  Matter \textbf{8}, 102263 (2025).

\bibitem{opencode2025}
  Anomaly Co.
  \textit{OpenCode AI --- open-source AI coding agent}.
  \url{https://github.com/anomalyco/opencode} (accessed 2026).

\bibitem{yao2023react}
  S.~Yao \textit{et al.}
  \textit{ReAct: Synergizing reasoning and acting in language models}. ICLR (2023).

\bibitem{schick2023toolformer}
  T.~Schick \textit{et al.}
  \textit{Toolformer: Language models can teach themselves to use tools}.
  NeurIPS (2023).

\bibitem{mcp2024}
  Anthropic.
  \textit{Model Context Protocol specification}.
  \url{https://modelcontextprotocol.io} (2024).

\bibitem{anthropic_claude}
  Anthropic.
  \textit{The Claude Model Family: Capabilities and Safety Evaluations}
  (Technical Reports, 2024--2025).

\bibitem{openai_gpt4}
  OpenAI.
  \textit{GPT-4 Technical Report}. arXiv:2303.08774 (2023).

\bibitem{ollama}
  Ollama.
  \textit{Ollama: Local large language model runtime}.
  \url{https://ollama.ai} (accessed 2026).

\bibitem{textual}
  Textualize.
  \textit{Textual: A TUI framework for Python}.
  \url{https://textual.textualize.io} (accessed 2026).

\bibitem{typer}
  S.~Ram\'irez.
  \textit{Typer: Build great CLIs}. \url{https://typer.tiangolo.com} (accessed 2026).

\bibitem{fastapi}
  S.~Ram\'irez.
  \textit{FastAPI: Modern, fast web framework for Python}.
  \url{https://fastapi.tiangolo.com} (accessed 2026).

\end{thebibliography}

% ============================================================
\vspace{2em}
\begin{center}
\small\textit{LRET Framework \textbullet\ Technical Report \textbullet\ April 2026
\textbullet\ Repository: \texttt{github.com/kunal5556/LRET}}
\end{center}
% ============================================================

\end{document}
